% !TeX spellcheck = pt_BR
\documentclass[12pt,portuguese]{article}

%----------------------------------------------------------------------
% PACOTES ESSENCIAIS E TRADUÇÃO
%----------------------------------------------------------------------
\usepackage[brazilian]{babel}     % Suporte à língua portuguesa
\usepackage[utf8]{inputenc}       % Codificação de entrada
\usepackage[T1]{fontenc}          % Codificação da fonte
\usepackage{lmodern}              % Fonte escalável

%----------------------------------------------------------------------
% PACOTES MATEMÁTICOS
%----------------------------------------------------------------------
\usepackage{amsmath, amsthm, amssymb, amstext, mathtools}
\usepackage{mathrsfs}             % Para fontes como \mathscr
\usepackage{dsfont}               % Para o indicador de conjunto \mathds{1}
\usepackage[all]{xy}              % Para diagramas comutativos
\usepackage{imakeidx}             % Criação de Lista de Símbolos / Índice

% Desigualdades sempre com símbolos inclinados
\let\leq\leqslant
\let\le\leqslant
\let\ge\geqslant
\let\geq\geqslant

%----------------------------------------------------------------------
% FORMATAÇÃO, LAYOUT E CORES
%----------------------------------------------------------------------
\usepackage{geometry}
\usepackage{varwidth}
\usepackage{graphicx}
\usepackage[pdftex]{xcolor}
\usepackage{titlesec}
\usepackage{appendix}
\usepackage{enumitem}
\usepackage{float}
\usepackage{booktabs}
\usepackage{array}
\usepackage{colortbl}
\usepackage{fontawesome5}
\usepackage{lucide-icons}
\usepackage{caption}
\usepackage{stackengine}
\captionsetup{width=0.8\textwidth}

% Definição de Cores
\definecolor{teal}{RGB}{15,118,110}
\definecolor{text}{RGB}{28,25,23}
\definecolor{muted}{RGB}{87,83,78}
\definecolor{card}{RGB}{255,255,255}
\definecolor{border}{RGB}{231,229,228}

% Definição de margens e espaçamento
\setlength{\topmargin}{-0.7in}
\setlength{\oddsidemargin}{0in}
\setlength{\textheight}{9.55in}
\setlength{\textwidth}{6.5in}
\setlength{\baselineskip}{12mm}

% Formatação dos títulos de seção
\titleformat{\section}{\Large\bfseries}{\thesection.}{0.5em}{}
\titleformat{\subsection}{\large\bfseries}{\thesubsection.}{0.5em}{}
\titleformat{\subsubsection}{\normalfont\itshape}{\thesubsubsection.}{0.6em}{}

%----------------------------------------------------------------------
% TIKZ E GRÁFICOS
%----------------------------------------------------------------------
\usepackage{tikz}
\usepackage{pgfplots}
\pgfplotsset{compat=1.17}
\usetikzlibrary{positioning,calc,fit,nfold,decorations.pathreplacing,calligraphy,tikzmark,patterns,arrows.meta,bending,hobby}

%----------------------------------------------------------------------
% REFERÊNCIAS, LINKS E METADADOS DO PDF
%----------------------------------------------------------------------
\usepackage[pdftex, colorlinks, linkcolor={red!50!black}, 
citecolor={red}, urlcolor={blue!80!black}]{hyperref}

\hypersetup{
pdftitle={Funções Características}, pdfauthor={L. Cioletti},
pdfsubject={Função característica, fórmula de inversão, teorema de
continuidade, momentos e derivadas}, pdfkeywords={função característica,
fórmula de inversão, convergência fraca, teorema de continuidade, momentos,
MAT-UnB} }

% --- PACOTE E COMANDO PARA A LICENÇA CREATIVE COMMONS ---
\usepackage[type={CC},modifier={by-nc-sa},version={4.0}]{doclicense}
\newcommand\blfootnote[1]{%
\begingroup
\renewcommand\thefootnote{}\footnote{#1}%
\addtocounter{footnote}{-1}%
\endgroup
}

%----------------------------------------------------------------------
% ALIASCNT E AMBIENTES DE TEOREMA (Cleveref)
%----------------------------------------------------------------------
\usepackage{aliascnt}
\theoremstyle{definition}
\newtheorem{theorem}{Teorema}

\newaliascnt{definition}{theorem}
\newtheorem{definition}[definition]{Definição}
\aliascntresetthe{definition}

\newaliascnt{example}{theorem}
\newtheorem{example}[example]{Exemplo}
\aliascntresetthe{example}

\newaliascnt{proposition}{theorem}
\newtheorem{proposition}[proposition]{Proposição}
\aliascntresetthe{proposition}

\newaliascnt{corollary}{theorem}
\newtheorem{corollary}[corollary]{Corolário}
\aliascntresetthe{corollary}

\newaliascnt{lemma}{theorem}
\newtheorem{lemma}[lemma]{Lema}
\aliascntresetthe{lemma}

\newaliascnt{remark}{theorem}
\newtheorem{remark}[remark]{Observação}
\aliascntresetthe{remark}

\newtheoremstyle{plainexercise}
  {3pt}{3pt}{\normalfont}{}{\bfseries}{.}{.5em}
  {\thmnumber{#2}\thmnote{ #3}}
\theoremstyle{plainexercise}
\newtheorem{exercise}{Exercício}

% Versão marcada dos exercícios: a estrela fica pendurada na margem esquerda,
% como nas notas anteriores.
\newtheoremstyle{starredexercise}
  {3pt}{3pt}{\normalfont}{}{\bfseries}{.}{.5em}
  {\llap{\textasteriskcentered\,}\thmnumber{#2}\thmnote{ #3}}
\theoremstyle{starredexercise}
\newtheorem{starredexercise}[exercise]{Exercício}
\theoremstyle{definition}

% Ambientes não numerados usados no texto
\newtheorem*{property}{Propriedade}
\newtheorem*{solution}{Solução}

% Se um enumerate for o primeiro conteúdo de um teorema/exercício (cabeçalho
% ainda pendente, \@inlabel true), o \leavevmode dispara o \everypar do amsthm
% que coloca o cabeçalho na própria linha, antes da lista iniciar.
\makeatletter
\setlist[enumerate]{before={\if@inlabel\leavevmode\fi}}
\makeatother

% Nomes em português para os ambientes
\AtBeginDocument{\renewcommand{\proofname}{\normalfont\textbf{Prova}}}
\renewcommand{\qedsymbol}{$\blacksquare$}

% Cleveref (deve vir depois de hyperref e amsthm)
\usepackage[brazilian,nameinlink]{cleveref}

\crefname{theorem}{teorema}{teoremas}
\Crefname{theorem}{Teorema}{Teoremas}
\crefname{definition}{definição}{definições}
\Crefname{definition}{Definição}{Definições}
\crefname{lemma}{lema}{lemas}
\Crefname{lemma}{Lema}{Lemas}
\crefname{proposition}{proposição}{proposições}
\Crefname{proposition}{Proposição}{Proposições}
\crefname{corollary}{corolário}{corolários}
\Crefname{corollary}{Corolário}{Corolários}
\crefname{example}{exemplo}{exemplos}
\Crefname{example}{Exemplo}{Exemplos}
\crefname{exercise}{exercício}{exercícios}
\Crefname{exercise}{Exercício}{Exercícios}
\crefname{remark}{observação}{observações}
\Crefname{remark}{Observação}{Observações}
\crefname{section}{seção}{seções}
\Crefname{section}{Seção}{Seções}

%----------------------------------------------------------------------
% MACROS ADICIONAIS
%----------------------------------------------------------------------
\let\oldemptyset\emptyset
\let\emptyset\varnothing
\graphicspath{{./}{../arquivos-aux/}}

% Seno na convenção em português
\renewcommand{\sin}{\operatorname{sen}}

%----------------------------------------------------------------------
% DOCUMENTO PRINCIPAL
%----------------------------------------------------------------------
\begin{document}

% --- CABEÇALHO COM O LOGO DA UnB ---
\vspace*{-0.8cm}
\begin{figure}[H]
\includegraphics{logo-mat-unb.pdf}
\end{figure}

% --- TÍTULO, AUTOR E DATA ---
\vspace*{0.3cm}
\begin{center}
{\LARGE\bf Funções Características}
\vspace{1cm}

{\large \bf L. Cioletti}
\vspace{0.3cm}

{\footnotesize Outubro de 2026}
\end{center}

% --- RESUMO (ABSTRACT) ---
\begin{abstract}
Nestas notas, introduzimos a função característica de uma variável aleatória e
mostramos que ela determina completamente a distribuição de $X$. Demonstramos
as propriedades elementares
das funções características, a fórmula de inversão e o teorema de
continuidade, que relaciona a convergência fraca de distribuições com a
convergência pontual das funções características correspondentes. Estudamos
também a relação entre o comportamento de $\varphi$ na origem e os momentos de
$X$, obtendo a expansão de segunda ordem e a recíproca que garante a
existência do segundo momento. Calculamos a função característica das
principais distribuições do curso e, por meio da fórmula de inversão,
obtemos as distribuições de Pólya e de Cauchy. Incluímos uma lista de
exercícios propostos. A exposição segue de perto a Seção 3.3 do Capítulo 3 de
\cite{MR3930614}.
\end{abstract}

% --- LICENÇA NO RODAPÉ ---
\blfootnote{%
\doclicenseImage[imagewidth=2.6em]\ 
\begingroup
\hypersetup{urlcolor=darkgray}
\href{\doclicenseURL}{%
\doclicenseName - Attribution-NonCommercial-ShareAlike 4.0 International%
}%
\endgroup
}

%----------------------------------------------------------------------
% CONTEÚDO
%----------------------------------------------------------------------
\section{Introdução}\label{sec:introducao}

Nas notas de aula anteriores sobre funções geradoras de momentos, estudamos a
função $M_X(t)=\mathbb{E}\bigl[e^{tX}\bigr]$ e vimos que ela gera os momentos
de $X$ por meio de suas derivadas na origem. O preço pago por essa comodidade é
uma hipótese de existência: a esperança $\mathbb{E}\bigl[e^{tX}\bigr]$ pode não
estar definida para nenhum $t\neq0$. O exemplo da distribuição de Cauchy, já
discutido naquelas notas, é o caso extremo, pois $M_X(t)=+\infty$ para todo
$t\neq0$. Para essa distribuição, a função geradora de momentos não apenas é
difícil de calcular: ela simplesmente não existe.

A função característica contorna essa dificuldade com uma modificação simples
no expoente. Trocamos o número real $t$ pelo número imaginário $it$, isto é,
passamos a considerar
\[
\varphi(t)=\mathbb{E}\bigl[e^{itX}\bigr].
\]
A vantagem é imediata. Pela fórmula de Euler,
$e^{itX}=\cos(tX)+i\sin(tX)$, e como
$\bigl|e^{itX}\bigr|=\cos^2(tX)+\sin^2(tX)=1$, a variável aleatória $e^{itX}$ é
limitada e, portanto, integrável. Logo $\varphi(t)$ está bem definida para toda
variável aleatória $X$ e para todo $t\in\mathbb{R}$, sem nenhuma hipótese sobre
os momentos de $X$. Esse ganho de generalidade tem um custo: $\varphi$ assume
valores complexos e as manipulações com ela exigem o cuidado de separar parte
real e parte imaginária.

O objetivo destas notas é desenvolver a teoria básica das funções
características em três frentes. Primeiro, mostramos que $\varphi$ determina a
distribuição de $X$ e apresentamos a fórmula de inversão, que recupera a medida
de probabilidade a partir de $\varphi$ por meio de uma integral na variável
$t$. Em seguida, relacionamos a convergência fraca de uma sequência de
distribuições com a convergência pontual das funções características
correspondentes, o que fornece uma das ferramentas mais eficientes para provar
teoremas de convergência. Por fim, estudamos a relação entre o comportamento de
$\varphi$ em uma vizinhança da origem e os momentos de $X$.

A organização é a seguinte. A \Cref{sec:preliminares} reúne os fatos auxiliares
sobre variáveis aleatórias complexas, derivação sob o sinal de integral e a
integral de Dirichlet, usados no restante do texto. A
\Cref{sec:definicao} apresenta a definição e as propriedades elementares. A
\Cref{sec:exemplos} calcula a função característica das principais
distribuições do curso. A \Cref{sec:inversao} é dedicada à fórmula de inversão e
a suas consequências. A \Cref{sec:continuidade} contém o teorema de
continuidade, e a \Cref{sec:momentos} trata da relação entre derivadas na
origem e momentos. Os exercícios propostos estão reunidos na
\Cref{sec:exercicios}.

Os resultados destas notas pressupõem os conceitos e as técnicas desenvolvidos
nas notas de aula anteriores sobre espaços de probabilidade, variáveis
aleatórias, esperança, funções de distribuição e modos de convergência. Em
particular, usaremos sem demonstrar a convergência dominada, a desigualdade de
Jensen e as caracterizações usuais da convergência fraca ali estabelecidas.

\section{Preliminares}\label{sec:preliminares}

Esta seção reúne três fatos técnicos que serão usados repetidamente. O leitor
pode consultá-la à medida que as referências apareçam no texto.

\subsection{Variáveis aleatórias complexas}\label{subsec:va-complexas}

Como $\varphi$ assume valores complexos, o primeiro passo é recordar o que
significa tomar a esperança de uma variável aleatória complexa.

\begin{definition}\label{def:va-complexa}
Dizemos que $Z:\Omega\to\mathbb{C}$ é uma \emph{variável aleatória complexa} se
$Z=X+iY$, onde $X$ e $Y$ são variáveis aleatórias reais. Nesse caso, definimos
\[
\mathbb{E}[Z] = \mathbb{E}[X]+i\,\mathbb{E}[Y],
\]
sempre que $\mathbb{E}\bigl[|X|\bigr]<\infty$ e $\mathbb{E}\bigl[|Y|\bigr]<\infty$.
\end{definition}

Escrevemos também $\operatorname{Re}Z=X$ e $\operatorname{Im}Z=Y$ para a parte
real e a parte imaginária de $Z$, e $\overline{Z}=X-iY$ para o conjugado.
Recorde que $|Z|=\sqrt{X^2+Y^2}$.

A definição é coerente porque as desigualdades
\[
|X|\leqslant |Z|\leqslant |X|+|Y|,
\qquad
|Y|\leqslant |Z|\leqslant |X|+|Y|
\]
mostram que $Z$ é integrável se, e somente se, suas partes real e imaginária o
são. Em particular, $e^{itX}$ é integrável para todo $t$, pois
$\bigl|e^{itX}\bigr|=1$. Temos ainda as propriedades usuais: a esperança é
linear sobre $\mathbb{C}$, comuta com a conjugação, isto é,
$\overline{\mathbb{E}[Z]}=\mathbb{E}\bigl[\overline{Z}\bigr]$, e satisfaz a
desigualdade triangular. Registramos esta última separadamente, pois ela será
usada na prova das propriedades das funções características.

\begin{lemma}\label{lem:esperanca-modulo}
Se $Z$ é uma variável aleatória complexa integrável, então
\[
\bigl|\mathbb{E}[Z]\bigr|\leqslant \mathbb{E}\bigl[|Z|\bigr].
\]
\end{lemma}

\begin{proof}
A função $\rho(x,y)=(x^2+y^2)^{1/2}$ é convexa em $\mathbb{R}^2$, pois é a
composição da norma euclidiana, que é convexa, com uma função crescente. Pela
desigualdade de Jensen em $\mathbb{R}^2$ aplicada ao par $(X,Y)$, obtemos
\[
\bigl|\mathbb{E}[Z]\bigr|
= \rho\bigl(\mathbb{E}[X],\mathbb{E}[Y]\bigr)
\leqslant \mathbb{E}\bigl[\rho(X,Y)\bigr]
= \mathbb{E}\bigl[|Z|\bigr],
\]
que é a desigualdade desejada.
\end{proof}

\subsection{Derivação sob o sinal de integral}\label{subsec:derivacao-integral}

Em vários pontos destas notas vamos derivar uma expressão da forma
$t\mapsto\int f(t,s)\,\nu(ds)$ e trocar a derivada com a integral. O resultado
a seguir é a versão que usaremos; ele é consequência direta do teorema do valor
médio e do teorema da convergência dominada.

\begin{theorem}[Derivação sob o sinal de integral]\label{thm:derivacao-integral}
Sejam $(S,\mathcal{S},\nu)$ um espaço de medida, $I\subseteq\mathbb{R}$ um
intervalo aberto e $f:I\times S\to\mathbb{C}$. Suponha que:
\begin{enumerate}
\item[(i)] para cada $t\in I$, a função $s\mapsto f(t,s)$ é integrável;
\item[(ii)] para $\nu$-quase todo $s$, a função $t\mapsto f(t,s)$ é derivável em
$I$;
\item[(iii)] existe $g:S\to[0,\infty]$ integrável tal que
$\bigl|\partial_t f(t,s)\bigr|\leqslant g(s)$ para todo $t\in I$ e
$\nu$-quase todo $s$.
\end{enumerate}
Então $F(t)=\int f(t,s)\,\nu(ds)$ é derivável em $I$ e
\[
F'(t)=\int \partial_t f(t,s)\,\nu(ds)
\qquad\text{para todo }t\in I.
\]
\end{theorem}

\begin{proof}
Basta tratar o caso real, aplicando o argumento separadamente às partes real e
imaginária de $f$. Fixe $t\in I$ e tome uma sequência $h_n\to0$ com $h_n\neq0$
e $t+h_n\in I$. Pelo teorema do valor médio, para $\nu$-quase todo $s$ existe
$\xi_n(s)$ entre $t$ e $t+h_n$ tal que
\[
\frac{f(t+h_n,s)-f(t,s)}{h_n}=\partial_t f(\xi_n(s),s).
\]
Por (iii), o quociente é majorado em módulo por $g(s)$, que não depende de $n$.
Além disso, por (ii), o quociente converge para $\partial_t f(t,s)$ quando
$n\to\infty$. A convergência dominada aplicada à medida $\nu$ fornece então
\[
\lim_{n\to\infty}\frac{F(t+h_n)-F(t)}{h_n}
= \int \partial_t f(t,s)\,\nu(ds).
\]
Como a sequência $(h_n)$ é arbitrária, o limite do quociente de diferenças
existe e não depende dela, o que conclui a prova.
\end{proof}

O papel da hipótese (iii) é fornecer uma majorante integrável que não dependa de
$t$. Sem ela, a troca entre derivada e integral pode falhar, mesmo que todas as
integrais envolvidas existam.

\subsection{A integral de Dirichlet}\label{subsec:dirichlet}

A prova da fórmula de inversão, na \Cref{sec:inversao}, depende do valor da
integral imprópria
\[
\int_0^\infty \frac{\sin x}{x}\,dx .
\]
Esse é um fato clássico de cálculo, mas como todos os seus passos serão usados
na prova do teorema, registramos o enunciado e a demonstração.

\begin{lemma}[Integral de Dirichlet]\label{lem:dirichlet}
\[
\lim_{T\to\infty}\int_0^T \frac{\sin x}{x}\,dx=\frac{\pi}{2}.
\]
\end{lemma}

\begin{proof}
O integrando é contínuo e limitado em $(0,\infty)$, com
$\lim_{x\to0}(\sin x)/x=1$. Dividimos a prova em três etapas.

\smallskip
\noindent\textbf{Etapa 1: a integral imprópria converge.}
Para $0<\varepsilon<T$, a integração por partes com $u=1/x$ e
$dv=\sin x\,dx$ fornece
\[
\int_\varepsilon^T \frac{\sin x}{x}\,dx
=
\left[\frac{1-\cos x}{x}\right]_\varepsilon^T
+
\int_\varepsilon^T \frac{1-\cos x}{x^2}\,dx.
\]
A função $(1-\cos x)/x^2$ é contínua e limitada em $(0,1]$, pois tende a $1/2$
quando $x\to0$, e é majorada por $2/x^2$ em $[1,\infty)$. Logo ela é integrável
em $(0,\infty)$ e a segunda parcela tem limite quando $\varepsilon\to0$ e
$T\to\infty$. O mesmo ocorre com a primeira, pois
$(1-\cos x)/x\leqslant 2/x$. Denotemos por $L$ o valor do limite.

\smallskip
\noindent\textbf{Etapa 2: uma representação auxiliar de $L$.}
Para $\varepsilon>0$, defina
\[
h(\varepsilon)=\int_0^\infty e^{-\varepsilon x}\,\frac{\sin x}{x}\,dx.
\]
A integral converge absolutamente, pois $|(\sin x)/x|\leqslant 1$ perto de $0$
e $e^{-\varepsilon x}/x$ é integrável em $[1,\infty)$. Fixado $\varepsilon_0>0$,
para todo $\varepsilon\geqslant\varepsilon_0$ vale
$\bigl|e^{-\varepsilon x}\sin x\bigr|\leqslant e^{-\varepsilon_0 x}$, que é
integrável em $(0,\infty)$. O \Cref{thm:derivacao-integral} permite então
derivar sob o sinal de integral e obter
\[
h'(\varepsilon)
=
-\int_0^\infty e^{-\varepsilon x}\sin x\,dx
=
-\frac{1}{1+\varepsilon^2},
\]
onde a última integral é elementar. Além disso,
$|h(\varepsilon)|\leqslant\int_0^\infty e^{-\varepsilon x}\,dx=1/\varepsilon$,
de modo que $h(\varepsilon)\to0$ quando $\varepsilon\to\infty$. Integrando
$h'$ e usando que $\arctan\varepsilon\to\pi/2$, concluímos que
\[
h(\varepsilon)=\frac{\pi}{2}-\arctan\varepsilon
\qquad\text{para todo }\varepsilon>0.
\]

\smallskip
\noindent\textbf{Etapa 3: passagem ao limite $\varepsilon\to0$.}
Seja $F(x)=\int_0^x(\sin u)/u\,du$ para $x>0$ e $F(x)=0$ para $x\leqslant0$.
Pela Etapa 1, $F$ é contínua, limitada, digamos por $C$, e $F(x)\to L$ quando
$x\to\infty$. A integração por partes em $(0,R)$ dá
\[
\int_0^R e^{-\varepsilon x}\,\frac{\sin x}{x}\,dx
=
e^{-\varepsilon R}F(R)+\varepsilon\int_0^R e^{-\varepsilon x}F(x)\,dx,
\]
e o termo de fronteira desaparece quando $R\to\infty$, pois
$e^{-\varepsilon R}F(R)\leqslant Ce^{-\varepsilon R}$. Logo
\[
h(\varepsilon)=\varepsilon\int_0^\infty e^{-\varepsilon x}F(x)\,dx.
\]
Fixe $\delta>0$ e escolha $X$ tal que $|F(x)-L|\leqslant\delta$ para
$x\geqslant X$. Escrevendo
\[
\varepsilon\int_0^\infty e^{-\varepsilon x}F(x)\,dx-L
=
\varepsilon\int_0^\infty e^{-\varepsilon x}\bigl(F(x)-L\bigr)\,dx
\]
e separando a integral em $(0,X)$ e $(X,\infty)$, obtemos
\[
\left|h(\varepsilon)-L\right|
\leqslant
2C\,\varepsilon\int_0^X e^{-\varepsilon x}\,dx+\delta
\leqslant
2C\,\varepsilon X+\delta.
\]
Fazendo $\varepsilon\to0$ e depois $\delta\to0$, segue que
$h(\varepsilon)\to L$. Comparando com a Etapa 2, concluímos que $L=\pi/2$.
\end{proof}

\section{A Definição e as Primeiras Propriedades}\label{sec:definicao}

Nesta seção, introduzimos o objeto central destas notas e estabelecemos as
propriedades que valem para toda função característica.

\subsection{A definição}\label{subsec:fc-definicao}

Na \Cref{def:va-complexa} fixamos o significado de $\mathbb{E}[Z]$ para uma
variável aleatória complexa $Z$, e no \Cref{lem:esperanca-modulo} vimos que o
módulo da esperança não excede a esperança do módulo. Com essas duas
ferramentas podemos definir o objeto central destas notas.

\begin{definition}[Função característica]\label{def:func-caracteristica}
Seja $X$ uma variável aleatória real. A \emph{função característica} de $X$ é a função
$\varphi\colon\mathbb{R}\to\mathbb{C}$ dada por
\[
  \varphi(t) \;=\; \mathbb{E}\bigl[e^{itX}\bigr],
  \qquad t\in\mathbb{R}.
\]
\end{definition}

Observe, antes de tudo, que a definição faz sentido para \emph{toda} variável
aleatória, sem qualquer hipótese sobre momentos. De fato, para cada $t$
fixado, a variável aleatória $e^{itX} = \cos tX + i\sin tX$ tem módulo
unitário:
\[
  \bigl|e^{itX}\bigr| \;=\; \sqrt{\cos^2 tX + \sin^2 tX} \;=\; 1 .
\]
Logo $e^{itX}$ é limitada e, portanto, integrável. Além disso, a \Cref{def:va-complexa}
nos diz que basta calcular as esperanças das partes real e imaginária,
que são as variáveis aleatórias limitadas $\cos tX$ e $\sin tX$.

A fórmula de Euler permite reescrever a definição de modo inteiramente real.
Como $e^{itX} = \cos tX + i\sin tX$, a definição de esperança de uma variável
aleatória complexa dá
\begin{equation}\label{eq:fc-partes}
  \varphi(t) \;=\; \mathbb{E}\bigl[\cos tX\bigr] + i\,\mathbb{E}\bigl[\sin tX\bigr],
  \qquad t\in\mathbb{R},
\end{equation}
e, separando as partes de $\varphi(t)$, obtemos
\[
  \operatorname{Re}\varphi(t) \;=\; \mathbb{E}\bigl[\cos tX\bigr]
  \qquad\text{e}\qquad
  \operatorname{Im}\varphi(t) \;=\; \mathbb{E}\bigl[\sin tX\bigr].
\]
Assim, $\operatorname{Re}\varphi$ e $\operatorname{Im}\varphi$ são funções
reais de variável real, e conhecer $\varphi$ equivale a conhecer esse par de
funções reais. Em particular, $\varphi(0) = \mathbb{E}[1] = 1$. Aplicando
ainda o \Cref{lem:esperanca-modulo} a $e^{itX}$ e usando que o seu módulo é
sempre igual a $1$, segue que
\[
  \bigl|\varphi(t)\bigr| \;\leqslant\; \mathbb{E}\bigl[\bigl|e^{itX}\bigr|\bigr] \;=\; 1
  \qquad\text{para todo } t\in\mathbb{R}.
\]

\medskip

Vale a pena registrar desde já o que se ganha e o que se perde com essa
definição, porque a comparação com a função geradora de momentos orienta toda
a exposição. O que se ganha é a universalidade: $\varphi$ existe para toda
variável aleatória $X$, sem que se exija que algum momento exista, e o seu
domínio é sempre a reta real inteira. Recorde que a função geradora de
momentos exige $\mathbb{E}[e^{sX}]<\infty$ para algum $s\neq 0$, condição que
pode falhar até para distribuições com todos os momentos finitos. O que se
perde é a simplicidade dos valores: $\varphi(t)$ é em geral um número
complexo, e não um número real não negativo, de modo que não podemos falar em
monotonicidade nem interpretar $\varphi(t)$ como uma massa ou um peso. A
informação sobre a distribuição de $X$ fica codificada em $\varphi$ de uma
forma menos direta, mas, como veremos, sem perda alguma.

\subsection{As propriedades básicas}\label{subsec:fc-propriedades}

As cinco propriedades de uma \Cref{def:func-caracteristica} reunidas abaixo
decorrem diretamente da definição e serão usadas em praticamente todos os
argumentos das próximas seções. Nesta ordem, elas registram a normalização na origem, a
simetria hermitiana, a limitação ao disco unitário, a regularidade da função e o
seu comportamento sob transformações afins.

\begin{theorem}[Propriedades básicas da função característica]\label{thm:propriedades-fc}
Seja $X$ uma variável aleatória real, com função característica $\varphi$. Para
todos $t,h,a,b\in\mathbb{R}$ valem:
\begin{enumerate}
\item[(a)] $\varphi(0)=1$;
\item[(b)] $\varphi(-t)=\overline{\varphi(t)}$;
\item[(c)] $|\varphi(t)|=\bigl|\mathbb{E}[e^{itX}]\bigr|\leqslant\mathbb{E}\bigl[|e^{itX}|\bigr]=1$;
\item[(d)] $|\varphi(t+h)-\varphi(t)|\leqslant\mathbb{E}\bigl[|e^{ihX}-1|\bigr]$; em particular, $\varphi$ é uniformemente contínua em $\mathbb{R}$;
\item[(e)] $\mathbb{E}\bigl[e^{it(aX+b)}\bigr]=e^{itb}\varphi(at)$.
\end{enumerate}
\end{theorem}

\begin{proof}
Os dois primeiros itens são consequências imediatas da definição e da estrutura
de $\mathbb{E}[\,\cdot\,]$ para variáveis aleatórias complexas, descrita em
\Cref{def:va-complexa}.

\smallskip
\noindent\textbf{Item (a).}
Como $e^{i0X}=1$, temos $\varphi(0)=\mathbb{E}[1]=1$. A função está, portanto,
normalizada na origem.

\smallskip
\noindent\textbf{Item (b).}
Escrevendo $\varphi$ em partes real e imaginária, como em
\Cref{def:va-complexa}, e usando a paridade do cosseno e a imparidade do seno,
\[
\varphi(-t)=\mathbb{E}\bigl[\cos(-tX)\bigr]+i\,\mathbb{E}\bigl[\sin(-tX)\bigr]
=\mathbb{E}\bigl[\cos(tX)\bigr]-i\,\mathbb{E}\bigl[\sin(tX)\bigr]
=\overline{\varphi(t)}.
\]
Equivalentemente, $e^{-itX}=\overline{e^{itX}}$ e a conjugação comuta com a
esperança, por atuar separadamente nas partes real e imaginária. A propriedade
diz que $\varphi$ é hermitiana; em particular, se $\varphi$ assume apenas
valores reais, então ela é uma função par.

\smallskip
\noindent\textbf{Item (c).}
A desigualdade é a do módulo da esperança, enunciada em
\Cref{lem:esperanca-modulo}: $|\mathbb{E}[Z]|\leqslant\mathbb{E}\bigl[|Z|\bigr]$
para toda variável aleatória complexa integrável $Z$. Esse lema é a desigualdade
de Jensen aplicada à função convexa $(x,y)\mapsto(x^2+y^2)^{1/2}$, que é a norma
euclidiana do plano, ao par $(\operatorname{Re}Z,\operatorname{Im}Z)$. Tomando
$Z=e^{itX}$ e lembrando que $|e^{itX}|=1$, obtemos
\[
|\varphi(t)|=\bigl|\mathbb{E}\bigl[e^{itX}\bigr]\bigr|
\leqslant\mathbb{E}\bigl[|e^{itX}|\bigr]=1,
\]
de modo que a imagem de $\varphi$ está contida no disco unitário fechado.

\smallskip
\noindent\textbf{Item (d).}
Pela linearidade da esperança e novamente por \Cref{lem:esperanca-modulo},
\[
\begin{aligned}
\bigl|\varphi(t+h)-\varphi(t)\bigr|
&=\bigl|\mathbb{E}\bigl[e^{i(t+h)X}-e^{itX}\bigr]\bigr|
\leqslant\mathbb{E}\bigl[\bigl|e^{i(t+h)X}-e^{itX}\bigr|\bigr]\\
&=\mathbb{E}\bigl[\bigl|e^{itX}\bigl(e^{ihX}-1\bigr)\bigr|\bigr]
=\mathbb{E}\bigl[|e^{ihX}-1|\bigr],
\end{aligned}
\]
pois $|e^{itX}|=1$. O ponto essencial é que o lado direito não depende de $t$.
Mostremos que ele tende a $0$ quando $h\to0$. O integrando
$Y_h\equiv|e^{ihX}-1|=2\bigl|\sin(hX/2)\bigr|$ é limitado por $2$ e tende a $0$
pontualmente, pela continuidade do seno e porque $hX(\omega)\to0$ para cada
$\omega$. Pelo teorema da convergência dominada, cujo majorante é a constante
$2$, segue que $\mathbb{E}\bigl[|e^{ihX}-1|\bigr]\to0$.

Como esse majorante independe de $t$, a mesma cota vale em toda a reta. Dado
$\varepsilon>0$, escolha $\delta>0$ tal que
$\mathbb{E}\bigl[|e^{ihX}-1|\bigr]<\varepsilon$ sempre que $|h|<\delta$; então
$|\varphi(t+h)-\varphi(t)|<\varepsilon$ para todo $t\in\mathbb{R}$ com
$|h|<\delta$. Isto é a continuidade uniforme de $\varphi$ em $\mathbb{R}$ e, em
particular, a continuidade em cada ponto.

\smallskip
\noindent\textbf{Item (e).}
Aqui basta usar as propriedades da exponencial complexa, a linearidade da
esperança e o fato de que $e^{itb}$ é constante, logo pode sair da esperança:
\[
\mathbb{E}\bigl[e^{it(aX+b)}\bigr]
=\mathbb{E}\bigl[e^{itb}\,e^{i(ta)X}\bigr]
=e^{itb}\,\mathbb{E}\bigl[e^{i(ta)X}\bigr]
=e^{itb}\,\varphi(at).
\]
Assim, se $Y\equiv aX+b$, a função característica de $Y$ é
$\varphi_Y(t)=e^{itb}\varphi_X(at)$. Essa é a regra de transformação afim da
função característica: uma translação multiplica $\varphi$ por uma exponencial
complexa de módulo $1$ e uma dilatação reescala o argumento. Vamos usá-la
repetidamente nos exemplos, para reduzir cada cálculo ao caso de uma variável
aleatória com parâmetros normalizados.
\end{proof}

\subsection{Somas de variáveis aleatórias independentes}\label{subsec:fc-soma}

A razão principal para introduzir funções características é o resultado a seguir:
ele converte somas de variáveis aleatórias independentes em produtos de funções.
Produtos são muito mais fáceis de analisar do que convoluções, de modo que,
sempre que a independência estiver presente, podemos trocar o problema analítico
de calcular uma convolução pelo problema algébrico de multiplicar funções
conhecidas.

\begin{theorem}[Função característica de uma soma de independentes]\label{thm:fc-soma-independentes}
Sejam $X_1$ e $X_2$ variáveis aleatórias independentes, com funções
características
\[
\varphi_j(t) = \mathbb{E}\big[e^{itX_j}\big], \qquad t \in \mathbb{R}, \quad j = 1,2.
\]
Então $X_1 + X_2$ tem função característica $\varphi_1(t)\varphi_2(t)$, isto é,
\[
\mathbb{E}\big[e^{it(X_1+X_2)}\big] = \varphi_1(t)\,\varphi_2(t),
\qquad t \in \mathbb{R}.
\]
\end{theorem}

\begin{proof}
Fixemos $t \in \mathbb{R}$. O argumento se organiza em três etapas.

\emph{(i) Fatoração da exponencial.} Pela propriedade básica da exponencial
complexa, $e^{a+b} = e^{a}e^{b}$ para quaisquer $a,b \in \mathbb{C}$. Tomando
$a = itX_1$ e $b = itX_2$, obtemos a identidade pontual
\[
e^{it(X_1+X_2)} = e^{itX_1}e^{itX_2},
\]
válida em todos os pontos do espaço de probabilidade.

\emph{(ii) Fatoração da esperança.} As variáveis aleatórias $e^{itX_1}$ e
$e^{itX_2}$ são integráveis, pois $\bigl|e^{itX_j}\bigr| = 1$ e portanto
$\mathbb{E}\bigl[\bigl|e^{itX_j}\bigr|\bigr] = 1 < \infty$. Elas também são
independentes: cada uma é uma função mensurável de $X_j$, e funções mensuráveis
de variáveis aleatórias independentes, consideradas em pares, continuam
independentes, como registramos ao estudar independência. Sendo assim, vale a
regra do produto para esperanças de variáveis independentes e integráveis,
\[
\mathbb{E}\big[e^{itX_1}e^{itX_2}\big]
= \mathbb{E}\big[e^{itX_1}\big]\,\mathbb{E}\big[e^{itX_2}\big].
\]

\emph{(iii) Conclusão.} Juntando as duas etapas e as definições de $\varphi_1$ e
$\varphi_2$,
\[
\mathbb{E}\big[e^{it(X_1+X_2)}\big]
= \mathbb{E}\big[e^{itX_1}e^{itX_2}\big]
= \mathbb{E}\big[e^{itX_1}\big]\,\mathbb{E}\big[e^{itX_2}\big]
= \varphi_1(t)\,\varphi_2(t).
\]
Como $t \in \mathbb{R}$ era arbitrário, isso é precisamente o enunciado.
\end{proof}

O resultado se estende sem esforço a qualquer número finito de parcelas.

\begin{corollary}\label{cor:fc-soma-n}
Se $X_1,\dots,X_n$ são variáveis aleatórias independentes, com funções
características $\varphi_1,\dots,\varphi_n$, então $X_1+\cdots+X_n$ tem função
característica $\prod_{j=1}^{n}\varphi_j(t)$.
\end{corollary}

\begin{proof}
Indução em $n$: o caso $n = 1$ é trivial e, supondo o resultado válido para
$n-1$ parcelas, aplicamos o \Cref{thm:fc-soma-independentes} à variável
$X_1+\cdots+X_{n-1}$, que é independente de $X_n$ (por ser função mensurável de
$X_1,\dots,X_{n-1}$), e obtemos
\[
\mathbb{E}\big[e^{it(X_1+\cdots+X_n)}\big]
= \Bigl(\prod_{j=1}^{n-1}\varphi_j(t)\Bigr)\varphi_n(t)
= \prod_{j=1}^{n}\varphi_j(t).
\]
\end{proof}

\section{Exemplos}\label{sec:exemplos}

Vamos agora calcular a função característica das distribuições mais usadas no
curso. Em todos os exemplos, o cálculo se reduz a uma integral ou a uma série
que envolve a exponencial complexa.

\subsection{Dois exemplos discretos}\label{subsec:exemplos-discretos}

Começamos pelos dois modelos discretos mais simples, a moeda equilibrada e a
distribuição de Poisson. Em ambos a função característica é uma soma de números
complexos, e a convergência dessa soma pode ser verificada diretamente.

\begin{example}[Moeda equilibrada]\label{ex:moeda}
Suponha que $X$ assuma apenas os valores $1$ e $-1$, com
$\mathbb{P}(X=1)=\mathbb{P}(X=-1)=\tfrac12$. Como
$\operatorname{Im}X=\{1,-1\}$ é finito, a esperança da variável aleatória
complexa $e^{itX}$ é uma soma finita. Pela definição de função característica e
pela fórmula para a esperança de uma função de uma variável aleatória discreta,
vista nas notas de aula sobre esperança, cada valor de $X$ entra com peso
$\tfrac12$:
\begin{align*}
\varphi(t)
&= \mathbb{E}\bigl[e^{itX}\bigr]
\\[0.3cm]
&= \sum_{x\in\operatorname{Im}X} e^{itx}\,\mathbb{P}(X=x)
\\[0.3cm]
&= \frac{e^{it}+e^{-it}}{2}.
\end{align*}
Pela fórmula de Euler, $e^{it}=\cos t+i\sin t$ e
$e^{-it}=\cos t-i\sin t$. Os senos se cancelam e os cossenos se somam, logo
\[
\varphi(t)=\cos t, \qquad t\in\mathbb{R}.
\]

Registramos três leituras do resultado. Primeiro, a conferência do item (a) de
\Cref{thm:propriedades-fc}: em $t=0$ obtemos $\varphi(0)=\cos 0=1$, como deve
ser. Segundo, $\varphi$ é real e par, o que reflete a simetria da distribuição:
trocar $X$ por $-X$ não altera a lei e corresponde a trocar $t$ por $-t$ dentro
da exponencial. Esse ponto será retomado em \Cref{cor:fc-real}, onde veremos que
a realidade da função característica é equivalente à simetria da distribuição.
Terceiro, $|\varphi(t)|=|\cos t|\leqslant 1$, em acordo com a cota geral
$|\varphi(t)|\leqslant\varphi(0)$.
\end{example}

\begin{example}[Distribuição de Poisson]\label{ex:poisson}
Seja $X\sim\mathrm{Poisson}(\lambda)$, com $\lambda>0$, isto é,
$\mathbb{P}(X=k)=e^{-\lambda}\lambda^k/k!$ para $k=0,1,2,\ldots$. Neste caso a
imagem é infinita, e a convergência da série se verifica pelos módulos. Como
$|e^{itk}|=1$ para todo $k$ e todo $t$, a série dos módulos não depende de $t$:
\begin{align*}
\sum_{k=0}^{\infty}
\Bigl|e^{-\lambda}\,\frac{\lambda^{k}e^{itk}}{k!}\Bigr|
&= e^{-\lambda}\sum_{k=0}^{\infty}\frac{\lambda^{k}}{k!}
\\[0.3cm]
&= e^{-\lambda}e^{\lambda}=1<\infty,
\end{align*}
onde usamos o desenvolvimento da exponencial. Logo a série converge absoluta e
uniformemente em $t$, o que autoriza agrupar e fatorar os termos livremente.

Escrevendo $e^{itk}=(e^{it})^{k}$, obtemos
\begin{align*}
\varphi(t)
&= \sum_{k=0}^{\infty} e^{itk}\,\mathbb{P}(X=k)
\\[0.3cm]
&= e^{-\lambda}\sum_{k=0}^{\infty}\frac{\bigl(\lambda e^{it}\bigr)^{k}}{k!}
\\[0.3cm]
&= e^{-\lambda}\exp\bigl(\lambda e^{it}\bigr)
\\[0.3cm]
&= \exp\bigl(\lambda(e^{it}-1)\bigr),
\end{align*}
pois $\sum_{k\geqslant 0}z^{k}/k!=e^{z}$ para todo $z\in\mathbb{C}$, com
$z=\lambda e^{it}$.

A conferência do item (a) de \Cref{thm:propriedades-fc} é imediata: em $t=0$
temos $\varphi(0)=\exp\bigl(\lambda(1-1)\bigr)=1$. Diferentemente do exemplo da
moeda, esta função característica não é real, pois
$\operatorname{Im}\varphi(t)=\exp(\lambda(\cos t-1))\sin(\lambda\sin t)$ não se
anula para todo $t$.
\end{example}

\subsection{A distribuição normal}\label{subsec:exemplo-normal}

\begin{example}[Distribuição normal padrão]\label{ex:normal}
Seja $Z$ uma variável aleatória com distribuição normal padrão, isto é, com
densidade $f(x)=(2\pi)^{-1/2}e^{-x^2/2}$, $x\in\mathbb{R}$. Vamos mostrar que a
função característica de $Z$ é
\[
\varphi(t)=\mathbb{E}\bigl[e^{itZ}\bigr]
=\int_{-\infty}^{\infty}e^{itx}\,(2\pi)^{-1/2}e^{-x^2/2}\,dx
=e^{-t^2/2}, \qquad t\in\mathbb{R}.
\]

Apresentamos dois argumentos para esse cálculo. O primeiro, que chamamos de
\emph{prova heurística}, completa um quadrado no expoente e desloca o contorno
de integração no plano complexo: ele identifica o valor correto supondo
legítima uma manipulação que não justificamos aqui, e por isso não é uma prova.
O segundo, a \emph{prova formal}, verifica o mesmo resultado por um caminho que
não depende daquela manipulação, ainda que informe pouco sobre a origem do
valor encontrado.
\end{example}

\begin{proof}[Prova heurística]
Começamos por reunir o integrando em um único exponencial,
$e^{itx}e^{-x^2/2}=e^{-(x^2/2-itx)}$, e completamos o quadrado no expoente.
Como $(x-it)^2=x^2-2itx-t^2$, dividindo por $2$ e reorganizando os termos
obtemos a identidade
\[
\frac{x^2}{2}-itx=\frac{(x-it)^2}{2}+\frac{t^2}{2}.
\]
Substituindo essa expressão na integral e usando a propriedade
$e^{a+b}=e^ae^b$ da exponencial, segue que
\begin{align*}
\varphi(t)
&=\int_{-\infty}^{\infty}(2\pi)^{-1/2}
\exp\biggl(-\frac{(x-it)^2}{2}-\frac{t^2}{2}\biggr)\,dx
\\[0.3cm]
&=e^{-t^2/2}\int_{-\infty}^{\infty}(2\pi)^{-1/2}
\exp\biggl(-\frac{(x-it)^2}{2}\biggr)\,dx ,
\end{align*}
pois o fator $e^{-t^2/2}$ não depende de $x$ e pode sair da integral.

A integral que resta é a integral sobre toda a reta da densidade de uma normal
com média $it$ e variância $1$. Se a média fosse um número real $a$, essa
integral valeria $1$, pela normalização da densidade normal. O argumento
consiste em estender essa identidade de $a$ real para o parâmetro complexo
$a=it$ por continuação analítica: com o deslocamento $x\mapsto x+it$, a
integração passa a ser feita ao longo da reta $\operatorname{Im}z=t$ e o valor
da integral é preservado. Esse é o passo que tomamos como legítimo sem
demonstrar, pois justificá-lo exigiria controlar o decaimento do integrando no
infinito e invocar a invariância da integral por deformação do contorno
(teorema de Cauchy). Por isso o argumento apenas identifica o valor correto da
função característica, e a verificação fica para a prova seguinte.
\end{proof}

\begin{proof}[Prova formal]
Primeiro usamos a paridade para eliminar a parte imaginária. Escrevendo
$e^{itx}=\cos(tx)+i\sin(tx)$, a função
$x\mapsto\sin(tx)(2\pi)^{-1/2}e^{-x^2/2}$ é ímpar e integrável em
$\mathbb{R}$, logo sua integral se anula e
\[
\varphi(t)=\int_{-\infty}^{\infty}\cos(tx)\,(2\pi)^{-1/2}e^{-x^2/2}\,dx .
\]
O integrando agora é real, e a majoração a seguir mostra que é legítimo derivar
sob o sinal de integral.

Para $t$ fixo, a derivada em $t$ do integrando é
$-x\sin(tx)(2\pi)^{-1/2}e^{-x^2/2}$, cujo módulo é dominado por
$g(x)=|x|(2\pi)^{-1/2}e^{-x^2/2}$. A função $g$ é integrável em
$\mathbb{R}$, pois $|x|e^{-x^2/2}$ decai exponencialmente, e a majorante não
depende de $t$. Pelo \Cref{thm:derivacao-integral}, podemos então derivar sob
o sinal de integral e obter
\[
\varphi'(t)=\int_{-\infty}^{\infty}\bigl(-x\sin(tx)\bigr)
(2\pi)^{-1/2}e^{-x^2/2}\,dx .
\]
Para calcular essa integral, observe que
\[
\frac{d}{dx}\Bigl[(2\pi)^{-1/2}e^{-x^2/2}\Bigr]
=-x\,(2\pi)^{-1/2}e^{-x^2/2},
\]
de modo que $-x\,(2\pi)^{-1/2}e^{-x^2/2}$ é a derivada da densidade normal
padrão e o integrando tem a forma $\sin(tx)f'(x)$. Integrando por partes em um
intervalo $[-R,R]$ e fazendo $R\to\infty$,
\begin{align*}
\varphi'(t)
&=\int_{-\infty}^{\infty}\sin(tx)\,\frac{d}{dx}
\Bigl[(2\pi)^{-1/2}e^{-x^2/2}\Bigr]\,dx
\\[0.3cm]
&=-\int_{-\infty}^{\infty}t\cos(tx)\,(2\pi)^{-1/2}e^{-x^2/2}\,dx
=-t\,\varphi(t),
\end{align*}
porque o termo de fronteira se anula, já que $e^{-x^2/2}\to 0$ quando
$|x|\to\infty$ enquanto $|\sin(tx)|\leqslant 1$, e a última integral é
exatamente a expressão de $\varphi(t)$ obtida acima.

Obtemos assim a equação diferencial linear de primeira ordem
$\varphi'(t)=-t\varphi(t)$, válida para todo $t\in\mathbb{R}$. Multiplicando
pelo fator integrante $e^{t^2/2}$,
\[
\frac{d}{dt}\Bigl[\varphi(t)e^{t^2/2}\Bigr]
=e^{t^2/2}\bigl(\varphi'(t)+t\varphi(t)\bigr)=0 ,
\]
logo $\varphi(t)e^{t^2/2}$ é constante em $\mathbb{R}$. Avaliando em $t=0$ e
lembrando que $\varphi(0)=\int_{-\infty}^{\infty}f(x)\,dx=1$, porque $f$ é
densidade, concluímos que essa constante vale $1$, isto é,
$\varphi(t)e^{t^2/2}=1$ para todo $t$ e, portanto,
$\varphi(t)=e^{-t^2/2}$, como queríamos.
\end{proof}

\medskip
Combinando o resultado com o item (e) de \Cref{thm:propriedades-fc}, obtemos a
função característica de qualquer membro da família normal. Se $X$ tem
distribuição normal com média $\mu$ e variância $\sigma^2>0$, então
$X$ tem a mesma distribuição que $\mu+\sigma Z$, com $Z$ normal padrão, e o item
(e) aplicado com $a=\sigma$ e $b=\mu$ fornece
\begin{align*}
\varphi_X(t)
&=\mathbb{E}\bigl[e^{it(\mu+\sigma Z)}\bigr]
=e^{i\mu t}\,\mathbb{E}\bigl[e^{i(\sigma t)Z}\bigr]
\\[0.3cm]
&=e^{i\mu t}\,\varphi_Z(\sigma t)
=\exp\Bigl(i\mu t-\frac{\sigma^2t^2}{2}\Bigr), \qquad t\in\mathbb{R}.
\end{align*}
Uma distribuição normal com média $\mu$ e variância $\sigma^2$ tem, portanto,
função característica $\exp\bigl(i\mu t-\sigma^2t^2/2\bigr)$.

O mesmo tipo de escalonamento se aplica aos outros exemplos, de modo que
frequentemente daremos a função característica apenas para um membro
conveniente da família, deixando os demais parâmetros para uma conta como a
acima.

\subsection{Três exemplos com densidade}\label{subsec:exemplos-continuos}

Nos três exemplos desta subseção a densidade é nula fora do intervalo
indicado. As integrais que definem a função característica se reduzem, por
isso, a integrais sobre intervalos limitados.

\begin{example}[Distribuição uniforme]\label{ex:uniforme}
Se $X$ tem densidade $f(x)=\frac{1}{b-a}\mathds{1}_{(a,b)}(x)$, com $-\infty<a<b<\infty$,
então
\[
\varphi(t)=\frac{e^{itb}-e^{ita}}{it\,(b-a)}, \qquad t\in\mathbb{R}\setminus\{0\},
\]
e $\varphi(0)=1$.
\end{example}

\begin{proof}
Para $\lambda\in\mathbb{C}$ não nulo vale $\frac{d}{dx}e^{\lambda x}=\lambda e^{\lambda x}$:
com $\lambda=\alpha+i\beta$, a fórmula de Euler escreve
$e^{\lambda x}=e^{\alpha x}\bigl(\cos(\beta x)+i\sin(\beta x)\bigr)$, e a regra do
produto aplicada às partes real e imaginária dá a identidade. Pelo teorema fundamental do
cálculo, aplicado a cada parte,
\[
\int_{a}^{b}e^{\lambda x}\,dx=\frac{e^{\lambda b}-e^{\lambda a}}{\lambda},
\]
com $\lambda$ no denominador porque a primitiva de $e^{\lambda x}$ é $e^{\lambda x}/\lambda$.
Tomando $\lambda=it$ com $t\neq0$ e dividindo pela massa total $b-a$ da densidade, obtemos
\[
\varphi(t)=\frac{1}{b-a}\int_{a}^{b}e^{itx}\,dx
=\frac{e^{itb}-e^{ita}}{it\,(b-a)}.
\]
Em $t=0$ temos $\varphi(0)=\int_{\mathbb{R}}f(x)\,dx=1$. A fórmula também tem limite $1$
quando $t\to0$, pois $e^{itx}=1+itx+O(t^{2}x^{2})$ fornece
$e^{itb}-e^{ita}=it(b-a)+O(t^{2})$ e o quociente por $it(b-a)$ tende a $1$. Logo a
expressão se estende por continuidade a $t=0$ com o valor $1$, forçado pela continuidade.
\end{proof}

No caso simétrico $a=-c$ e $b=c$, com $c>0$, o numerador é $e^{itc}-e^{-itc}=2i\sin(ct)$
e o denominador é $2ict$, de modo que $\varphi(t)=\sin(ct)/(ct)$ para $t\neq0$, com valor
$1$ em $t=0$; para $c=1$ isso dá $\sin t/t$. Vale comparar com \Cref{ex:moeda}: lá a
variável toma apenas os valores $-1$ e $1$ com probabilidade $1/2$ e $\varphi(t)=\cos t$,
média de $e^{itx}$ sobre esses dois pontos; aqui a média é tomada sobre todo o intervalo
$(-c,c)$, o que produz $\sin(ct)/(ct)$. O caso uniforme é, nesse sentido, a versão contínua
do exemplo da moeda.

\medskip

\begin{example}[Distribuição triangular]\label{ex:triangular}
Se $X$ tem densidade $f(x)=(1-|x|)\mathds{1}_{(-1,1)}(x)$, então
\[
\varphi(t)=\frac{2\,(1-\cos t)}{t^{2}}, \qquad t\neq0,
\]
com $\varphi(0)=1$. O nome vem do gráfico de $f$, que sobe de $0$ a $1$ e desce.
\end{example}

\begin{proof}
A densidade triangular é a da soma de duas variáveis uniformes independentes.
Sejam $U$ e $V$ independentes, ambas com densidade constante igual a $1$ em
$(-1/2,1/2)$ e nula fora dele. Pela fórmula da convolução,
\[
f_{U+V}(s)=\int_{-1/2}^{1/2}\mathds{1}_{(-1/2,1/2)}(s-y)\,dy,
\]
porque a densidade de $V$ vale $1$ em todo o intervalo de integração. O
integrando é $1$ exatamente quando $y\in(-1/2,1/2)$ e $s-y\in(-1/2,1/2)$, isto é,
quando $y\in(s-1/2,s+1/2)$; a integral é, portanto, o comprimento da interseção
$(-1/2,1/2)\cap(s-1/2,s+1/2)$. Para $s\in[0,1)$ essa interseção é
$(s-1/2,1/2)$, de comprimento $1-s$; para $s\in(-1,0]$ o cálculo simétrico dá
$1+s$; e para $|s|\geqslant1$ ela é vazia. Logo
$f_{U+V}(s)=(1-|s|)\mathds{1}_{(-1,1)}(s)$, que é precisamente a densidade de $X$.
Como a função característica depende apenas da distribuição, basta calcular $\varphi_{U+V}$.

Por \Cref{ex:uniforme}, com $a=-1/2$ e $b=1/2$, cada uma das variáveis tem
\[
\varphi_{U}(t)=\frac{e^{it/2}-e^{-it/2}}{it}=\frac{2\sin(t/2)}{t}, \qquad t\neq0,
\]
e $\varphi_{U}(0)=1$. Pela independência, \Cref{thm:fc-soma-independentes} diz
que a função característica da soma é o produto das funções características:
\[
\varphi_{X}(t)=\left(\frac{2\sin(t/2)}{t}\right)^{2}, \qquad t\neq0.
\]
A identidade $\cos 2\theta=1-2\sin^{2}\theta$ com $\theta=t/2$ fornece
$\sin^{2}(t/2)=\bigl(1-\cos t\bigr)/2$; substituindo no quadrado,
$\varphi_{X}(t)=4(1-\cos t)/(2t^{2})=2(1-\cos t)/t^{2}$.

Em $t=0$ a primeira forma tende a $1$, pois $\sin(t/2)\sim t/2$, e a segunda
também, pois $1-\cos t\sim t^{2}/2$. Os dois limites coincidem com
$\varphi_{X}(0)=\int_{\mathbb{R}}f(x)\,dx=1$, de modo que as expressões
obtidas para $t\neq0$ se estendem continuamente à origem. A função é ainda
real e par, como era de se esperar de uma densidade simétrica em torno de $0$.
\end{proof}

\medskip

\begin{example}[Distribuição exponencial]\label{ex:exponencial}
Se $X$ tem densidade $f(x)=e^{-x}\mathds{1}_{(0,\infty)}(x)$, então
\[
\varphi(t)=\frac{1}{1-it}, \qquad t\in\mathbb{R}.
\]
\end{example}

\begin{proof}
Por definição,
$\varphi(t)=\int_{0}^{\infty}e^{itx}e^{-x}\,dx=\int_{0}^{\infty}e^{(it-1)x}\,dx$.
A integral converge absolutamente, pois $|e^{itx}e^{-x}|=e^{-x}$ é integrável
em $(0,\infty)$; logo podemos calculá-la como limite de integrais em $[0,R]$.
Como $it-1\neq0$ para todo $t$ real, a primitiva é $e^{(it-1)x}/(it-1)$ e
\[
\int_{0}^{R}e^{(it-1)x}\,dx=\frac{e^{(it-1)R}-1}{it-1}.
\]
O termo de fronteira superior se anula, já que
$\bigl|e^{(it-1)R}\bigr|=\bigl|e^{itR}\bigr|e^{-R}=e^{-R}\to0$ quando
$R\to\infty$: o fator oscilante tem módulo $1$ e o decaimento exponencial domina.
Passando ao limite,
\[
\varphi(t)=\frac{0-1}{it-1}=\frac{1}{1-it},
\]
fórmula válida para todo $t$ real, sem extensão por continuidade, porque $1-it$
nunca se anula. Em particular $\varphi(0)=1$.
\end{proof}

\subsection{Misturas convexas}\label{subsec:misturas}

Vamos registrar uma operação simples sobre distribuições que preserva a boa
comportabilidade da função característica. Sejam $F_1,\dots,F_n$ funções de
distribuição e $\lambda_1,\dots,\lambda_n\geqslant0$ com soma $1$. A função
\[
  F(x) \;=\; \sum_{i=1}^{n}\lambda_i F_i(x), \qquad x\in\mathbb{R},
\]
é novamente uma função de distribuição, e $F$, ou a medida
$\mu=\sum_i\lambda_i\mu_i$ a ela associada, é a \emph{mistura convexa} de
$F_1,\dots,F_n$ com pesos $\lambda_1,\dots,\lambda_n$. Em termos
probabilísticos, sorteamos um índice $I$ com $\mathbb{P}(I=i)=\lambda_i$ e
observamos uma variável cuja lei, dado $\{I=i\}$, é $F_i$.

O resultado a seguir diz que a função característica da mistura é a mistura
das funções características com os mesmos pesos, como era de esperar: a
função característica é a integral de $x\mapsto e^{itx}$ contra a medida e a
integral é linear na medida, de modo que ao misturar as medidas as integrais
se misturam com os mesmos pesos.

\begin{lemma}[Misturas convexas]\label{lem:mistura}
Sejam $F_1,\dots,F_n$ funções de distribuição com funções características
$\varphi_1,\dots,\varphi_n$ e sejam $\lambda_1,\dots,\lambda_n\geqslant0$ com
$\lambda_1+\cdots+\lambda_n=1$. Então $\sum_{i=1}^{n}\lambda_iF_i$ tem função
característica $\sum_{i=1}^{n}\lambda_i\varphi_i$.
\end{lemma}

\begin{proof}
Seja $\mu=\sum_{i=1}^{n}\lambda_i\mu_i$ a medida da mistura, de modo que
$F(x)=\mu((-\infty,x])$. A integral contra uma soma de medidas se decompõe
como a soma das integrais, $\int f\,d(\mu+\nu)=\int f\,d\mu+\int f\,d\nu$, e
o mesmo vale para múltiplos não negativos; iterando para os $n$ termos da
mistura, obtemos, para todo $t\in\mathbb{R}$,
\[
  \varphi(t) = \int e^{itx}\,d\mu(x)
  = \sum_{i=1}^{n}\lambda_i\int e^{itx}\,d\mu_i(x)
  = \sum_{i=1}^{n}\lambda_i\,\varphi_i(t).
\]
\end{proof}

Em particular, o conjunto das funções características é convexo. O exemplo
seguinte mostra que essa observação elementar já produz uma função
característica pouco óbvia a partir da densidade.

\begin{example}[Exponencial bilateral]\label{ex:exponencial-bilateral}
A distribuição com densidade $f(x)=\tfrac12e^{-|x|}$, $x\in\mathbb{R}$,
conhecida como \emph{exponencial bilateral} (ou distribuição de Laplace), tem
função característica
\[
  \varphi(t) \;=\; \frac{1}{1+t^2}, \qquad t\in\mathbb{R}.
\]
\end{example}

\begin{proof}
Seja $X$ exponencial de parâmetro $1$, com densidade $e^{-x}$ em
$(0,\infty)$, e sejam $F_1$ a lei de $X$ e $F_2$ a lei de $-X$. Com
$\lambda_1=\lambda_2=\tfrac12$, a mistura $\tfrac12F_1+\tfrac12F_2$ tem
densidade $\tfrac12e^{-|x|}$: para $x>0$ ela vale $\tfrac12e^{-x}$ e, para
$x<0$, $\tfrac12e^{x}$, pois a densidade de $-X$ em $x$ é a de $X$ em $-x$.
Logo podemos aplicar \Cref{lem:mistura}. Por \Cref{ex:exponencial},
$\varphi_1(t)=1/(1-it)$ e, pelo item (b) de \Cref{thm:propriedades-fc}, a
função característica de $-X$ se obtém trocando $t$ por $-t$, de modo que
$\varphi_2(t)=\varphi_1(-t)=1/(1+it)$. Assim,
\[
  \varphi(t) = \tfrac12\varphi_1(t)+\tfrac12\varphi_2(t)
  = \frac{1}{2(1-it)}+\frac{1}{2(1+it)}
  = \frac{(1+it)+(1-it)}{2(1-it)(1+it)}
  = \frac{1}{1+t^2}.
\]
\end{proof}

Vale confrontar o resultado com a simetria da densidade: como $f(-x)=f(x)$,
a distribuição é simétrica em torno da origem e a função característica
obtida é real e par. De fato, as partes imaginárias vindas de $x$ e de $-x$
se cancelam duas a duas, restando apenas $\int\cos(tx)\,f(x)\,dx$.

\section{A Fórmula de Inversão}\label{sec:inversao}

Nos exemplos da seção anterior, partimos da distribuição e calculamos a função
característica. A pergunta natural é se esse procedimento pode ser invertido,
isto é, se a função característica determina a distribuição e, em caso
afirmativo, como recuperá-la. A resposta é afirmativa, e a fórmula que faz a
recuperação é o conteúdo desta seção.

\subsection{O teorema de inversão}\label{subsec:teorema-inversao}

A questão que se coloca agora é se a função característica determina a
distribuição. O resultado a seguir responde afirmativamente e fornece uma
fórmula que recupera a massa de intervalos a partir de $\varphi$; dele
dependem o critério de continuidade e a fórmula da densidade, adiante.

\begin{theorem}[Fórmula de inversão]\label{thm:inversao}
Seja $\mu$ uma medida de probabilidade em $\mathbb{R}$ e seja
$\varphi(t)=\int_{\mathbb{R}}e^{itx}\,\mu(dx)$ sua função característica. Se
$a<b$, então
\begin{equation}\label{eq:inversao}
\lim_{T\to\infty}\frac{1}{2\pi}\int_{-T}^{T}\frac{e^{-ita}-e^{-itb}}{it}\,\varphi(t)\,dt
=\mu\bigl((a,b)\bigr)+\tfrac12\,\mu\bigl(\{a,b\}\bigr).
\end{equation}
\end{theorem}

Escrevemos $\mu(a,b)$ como abreviação de $\mu\bigl((a,b)\bigr)$. O lado
direito de \eqref{eq:inversao} atribui a cada extremo metade de sua massa;
quando $\mu$ não tem átomos em $a$ nem em $b$, ele se reduz a
$\mu\bigl((a,b)\bigr)$.

\begin{remark}[A existência do limite faz parte da conclusão]\label{rem:inversao-limite}
O enunciado não supõe que o limite exista: a existência é parte da conclusão.
Isso merece atenção porque o integrando de \eqref{eq:inversao} não é, em
geral, absolutamente integrável em $\mathbb{R}$, de modo que o limite só tem
sentido como limite de truncamentos \emph{simétricos}, por cancelamento.

Para ver isso no caso mais simples, tome $\mu=\delta_0$, a massa unitária em $0$. Então
$\varphi(t)\equiv1$ e, com $a=-1$ e $b=1$, o integrando de \eqref{eq:inversao} vale
\[
\frac{e^{it}-e^{-it}}{it}=\frac{2i\sin t}{it}=\frac{2\sin t}{t},
\]
que decai como $1/t$ e cujo módulo não é integrável: em $[k\pi,(k+1)\pi]$
temos $\int_{k\pi}^{(k+1)\pi}|\sin t|\,dt=2$ e
$1/t\geqslant1/\bigl((k+1)\pi\bigr)$, donde
$\int_1^\infty\frac{|\sin t|}{t}\,dt\geqslant\frac{2}{\pi}\sum_{k\geqslant1}\frac{1}{k+1}=\infty$,
pela divergência da série harmônica. A integral imprópria converge apenas pela
compensação entre as partes positiva e negativa, e é ela que exige o corte
simétrico em $[-T,T]$. A fórmula, porém, vale: aqui
$S(T)=\int_0^T\frac{\sin t}{t}\,dt$ satisfaz
$\int_{-T}^{T}\frac{2\sin t}{t}\,dt=4S(T)\to2\pi$, e o lado esquerdo de
\eqref{eq:inversao} tende a $1$, o valor de
$\mu\bigl((-1,1)\bigr)+\tfrac12\mu\bigl(\{-1,1\}\bigr)$.
\end{remark}

\begin{proof}
A prova usa três ingredientes: uma identidade que controla o integrando perto
de $t=0$ e autoriza trocar a ordem das integrais, uma redução a funções do
tipo $\sin(\theta t)/t$ e o limite de Dirichlet do
\Cref{lem:dirichlet}. Indicamos o papel de cada etapa.

\smallskip
\noindent\textbf{Etapa 1: a expressão como integral dupla.}
Denotamos por $I_T$ a integral de \eqref{eq:inversao} sem o fator $1/2\pi$.
Substituindo a definição de $\varphi$,
\[
I_T=\int_{-T}^{T}\frac{e^{-ita}-e^{-itb}}{it}\,\varphi(t)\,dt
=\int_{-T}^{T}\int_{\mathbb{R}}\frac{e^{-ita}-e^{-itb}}{it}\,e^{itx}\,\mu(dx)\,dt ,
\]
e a meta é integrar primeiro em $t$, para que o integrando externo seja uma
função de $x$ cujo limite possamos calcular ponto a ponto.

\smallskip
\noindent\textbf{Etapa 2: o integrando é limitado, logo vale Fubini.}
O denominador $it$ sugere um problema em $t=0$, mas a dificuldade é apenas
aparente: para $t\neq0$, ao integrar a exponencial complexa e usar
$|e^{-ity}|=1$, obtemos
\[
\frac{e^{-ita}-e^{-itb}}{it}=\int_{a}^{b}e^{-ity}\,dy
\qquad\text{e}\qquad
\Bigl|\frac{e^{-ita}-e^{-itb}}{it}\Bigr|\leqslant\int_{a}^{b}dy=b-a .
\]
A identidade mostra também que o quociente tem limite $b-a$ quando $t\to0$, de
modo que o integrando se estende continuamente à origem e a cota vale em toda
a reta.

Essa é exatamente a informação que legitima o teorema de Fubini: o produto
$[-T,T]\times\mathbb{R}$ tem medida finita para $dt\otimes\mu$, pois $T$ é
finito e $\mu$ é medida de probabilidade, e o integrando, multiplicado por
$|e^{itx}|=1$, é limitado por $b-a$ nesse produto. Logo a integral dupla é
absolutamente convergente e as duas ordens de integração coincidem.

\smallskip
\noindent\textbf{Etapa 3: redução a senos.}
Integrando primeiro em $t$, o numerador se reescreve como
$\frac{e^{it(x-a)}-e^{it(x-b)}}{it}$. Para $\theta$ fixo, a fórmula de Euler e
as paridades $\cos(-\theta t)=\cos(\theta t)$ e
$\sin(-\theta t)=-\sin(\theta t)$ dão
\[
\frac{e^{i\theta t}}{it}=\frac{\cos(\theta t)}{it}+\frac{\sin(\theta t)}{t},
\]
e o primeiro termo é ímpar em $t$, pois o cosseno é par e $1/t$ é ímpar. As
contribuições desse termo sobre $[-T,0)$ e sobre $(0,T]$ se cancelam
exatamente, e a decomposição é legítima porque o integrando completo é
limitado: ele é igual a $\int_0^\theta e^{ity}\,dy$. Sobra
$\sin(\theta t)/t$, que é par, e assim
$\int_{-T}^{T}\frac{e^{i\theta t}}{it}\,dt=\int_{-T}^{T}\frac{\sin(\theta t)}{t}\,dt$;
tomando $\theta=x-a$ e $\theta=x-b$, obtemos
\begin{equation}\label{eq:IT}
I_T=\int_{\mathbb{R}}\Bigl\{
\int_{-T}^{T}\frac{\sin\bigl(t(x-a)\bigr)}{t}\,dt
-\int_{-T}^{T}\frac{\sin\bigl(t(x-b)\bigr)}{t}\,dt
\Bigr\}\mu(dx).
\end{equation}

\smallskip
\noindent\textbf{Etapa 4: isolar a dependência em $\theta$.}
Em \eqref{eq:IT} a variável $x$ aparece somente através das diferenças $x-a$ e
$x-b$. Definimos, para $\theta\in\mathbb{R}$ e $T>0$,
\[
R(\theta,T)\equiv\int_{-T}^{T}\frac{\sin(\theta t)}{t}\,dt,
\qquad\text{donde}\qquad
I_T=\int_{\mathbb{R}}\bigl\{R(x-a,T)-R(x-b,T)\bigr\}\,\mu(dx).
\]

\smallskip
\noindent\textbf{Etapa 5: relação com o truncamento de Dirichlet.}
Seja $S(T)\equiv\int_{0}^{T}\frac{\sin x}{x}\,dx$, o truncamento da integral
de Dirichlet. Para $\theta>0$ o integrando de $R(\theta,T)$ é par em $t$,
donde $R(\theta,T)=2\int_{0}^{T}\frac{\sin(\theta t)}{t}\,dt$, e a mudança de
variável $t=x/\theta$, com $dt=dx/\theta$, fornece
\[
R(\theta,T)=2\int_{0}^{T\theta}\frac{\sin x}{x/\theta}\,\frac{dx}{\theta}
=2\int_{0}^{T\theta}\frac{\sin x}{x}\,dx=2S(T\theta).
\]
Para $\theta<0$, a relação $\sin(\theta t)=-\sin\bigl(|\theta|t\bigr)$ dá
$R(\theta,T)=-R\bigl(|\theta|,T\bigr)$. Introduzindo a função sinal, com
$\operatorname{sgn}\theta=1$ se $\theta>0$, $\operatorname{sgn}\theta=-1$ se
$\theta<0$ e $\operatorname{sgn}\theta=0$ se $\theta=0$, as duas fórmulas se
reúnem em
\begin{equation}\label{eq:RT}
R(\theta,T)=2\,(\operatorname{sgn}\theta)\,S\bigl(T|\theta|\bigr),
\qquad \theta\in\mathbb{R},
\end{equation}
igualdade que também vale para $\theta=0$, caso em que ambos os lados são nulos.

\smallskip
\noindent\textbf{Etapa 6: limite pontual.}
Pelo \Cref{lem:dirichlet}, $S(T)\to\frac{\pi}{2}$ quando $T\to\infty$. Se
$\theta\neq0$, então $T|\theta|\to\infty$ e \eqref{eq:RT} dá
$R(\theta,T)\to\pi\operatorname{sgn}\theta$; para $\theta=0$ temos $R(0,T)=0$,
de modo que a conclusão vale em toda a reta. Com $\theta=x-a$ e $\theta=x-b$,
\[
\lim_{T\to\infty}\bigl\{R(x-a,T)-R(x-b,T)\bigr\}
=\begin{cases}
2\pi, & a<x<b,\\[0.2cm]
\pi, & x=a\ \text{ou}\ x=b,\\[0.2cm]
0, & x<a\ \text{ou}\ x>b.
\end{cases}
\]
De fato, se $a<x<b$ os sinais são opostos, $\operatorname{sgn}(x-a)=1$ e
$\operatorname{sgn}(x-b)=-1$, e a diferença vale $2\pi$; se $x$ é um dos
extremos, um sinal se anula e o outro vale $\pm1$, donde $\pi$; se $x$ está
fora de $[a,b]$, os sinais coincidem e a diferença é nula.

\smallskip
\noindent\textbf{Etapa 7: passagem do limite para dentro da integral.}
A cota que permite trocar limite e integral vem de \eqref{eq:RT} e não depende de $x$: como
$|\operatorname{sgn}\theta|\leqslant1$,
\[
|R(\theta,T)|\leqslant2\sup_{y\geqslant0}|S(y)|\equiv2M<\infty,
\qquad\text{logo}\qquad
\bigl|R(x-a,T)-R(x-b,T)\bigr|\leqslant4M .
\]
Aqui $M$ é finito porque $S$ é contínua em $[0,\infty)$ (o quociente
$\sin x/x$ se estende continuamente a $x=0$ com valor $1$ ) e tem limite
finito $\frac{\pi}{2}$ no infinito, e uma função contínua em $[0,\infty)$ com
limite finito no infinito é limitada. A constante $4M$ é integrável com
relação a $\mu$, que é medida de probabilidade.

Podemos então aplicar o teorema da convergência dominada. Ao longo de uma sucessão arbitrária
$T_n\to\infty$,
\[
I_{T_n}\longrightarrow
\pi\int_{\mathbb{R}}\bigl\{\operatorname{sgn}(x-a)-\operatorname{sgn}(x-b)\bigr\}\,\mu(dx),
\]
e, como o valor não depende da sucessão escolhida, o critério de subsequências
para convergência de números reais garante a convergência de $I_T$. Resta
identificar a integral. Como
\[
\tfrac12\bigl\{\operatorname{sgn}(x-a)-\operatorname{sgn}(x-b)\bigr\}
=\begin{cases}
1, & a<x<b,\\[0.2cm]
\tfrac12, & x=a\ \text{ou}\ x=b,\\[0.2cm]
0, & x<a\ \text{ou}\ x>b,
\end{cases}
\]
a função vale $1$ no intervalo aberto, $\tfrac12$ em cada extremo e $0$ fora
de $[a,b]$. Integrando contra $\mu$ e dividindo a expressão anterior por
$2\pi$,
\[
\frac{1}{2\pi}\,I_T\longrightarrow
\int_{\mathbb{R}}\tfrac12\bigl\{\operatorname{sgn}(x-a)-\operatorname{sgn}(x-b)\bigr\}\,\mu(dx)
=\mu\bigl((a,b)\bigr)+\tfrac12\,\mu\bigl(\{a,b\}\bigr),
\]
que é exatamente \eqref{eq:inversao}.
\end{proof}

\medskip
O \Cref{thm:inversao} mostra que $\varphi$ determina $\mu$ por completo. De
fato, escolhendo $a$ e $b$ fora do conjunto enumerável dos átomos de $\mu$, o
lado direito de \eqref{eq:inversao} é apenas $\mu\bigl((a,b)\bigr)$, de modo
que $\varphi$ determina a massa de todo intervalo aberto com extremos fora de
um conjunto enumerável; como os intervalos geram a $\sigma$ -álgebra de Borel
e duas medidas de probabilidade que coincidem em um $\pi$ -sistema gerador
coincidem em toda a $\sigma$ -álgebra, $\mu$ fica univocamente determinada. Em
particular, duas variáveis aleatórias com a mesma função característica têm a
mesma distribuição.

Vale também ler \eqref{eq:inversao} como uma média ponderada: o lado esquerdo
integra $\varphi(t)$ contra o peso $\frac{e^{-ita}-e^{-itb}}{2\pi i t}$, que
depende apenas do intervalo $(a,b)$, e o resultado é a massa que $\mu$ atribui
a esse intervalo. O fator $1/t$ explica por que o limite em $T$ é delicado: o
decaimento lento do peso impede integrabilidade absoluta e faz da
integrabilidade de $\varphi$ uma hipótese genuinamente restritiva, sob a qual
obteremos, na próxima subseção, a densidade associada a $\varphi$.

\subsection{Duas consequências imediatas}\label{subsec:consequencias-inversao}

A fórmula de inversão tem duas consequências imediatas, que provamos a seguir:
uma condição de simetria em termos de $\varphi_X$ e a estabilidade das normais
centradas sob convolução.

\begin{corollary}[Função característica real]\label{cor:fc-real}
Se $\varphi_X$ assume apenas valores reais, então $X$ e $-X$ têm a mesma distribuição.
\end{corollary}

\begin{proof}
Calculemos a função característica de $-X$. Por definição,
\[
  \varphi_{-X}(t) = \mathbb{E}\bigl[e^{it(-X)}\bigr]
  = \mathbb{E}\bigl[e^{-itX}\bigr], \qquad t \in \mathbb{R}.
\]
Como $e^{-i\theta} = \overline{e^{i\theta}}$ para todo $\theta$ real,
o integrando é o conjugado de $e^{itX}$ ponto a ponto; tomando esperanças
(calculadas parte real e parte imaginária a parte) obtemos
\[
  \varphi_{-X}(t) = \overline{\mathbb{E}\bigl[e^{itX}\bigr]}
  = \overline{\varphi_X(t)},
\]
que é o item (b) de \Cref{thm:propriedades-fc}. Por hipótese
$\varphi_X(t)$ é real, logo vale
$\varphi_{-X}(t) = \varphi_X(t)$ para todo $t$; pela unicidade dada pelo
\Cref{thm:inversao}, duas variáveis com a mesma função característica têm
a mesma distribuição, de modo que $X \overset{d}{=} -X$.
\end{proof}

\begin{corollary}[Soma de normais centradas independentes]\label{cor:soma-normais}
Sejam $X_1$ e $X_2$ independentes, com $X_i$ normal de média $0$ e variância
$\sigma_i^2$, $i = 1, 2$. Então $X_1 + X_2$ é normal de média $0$ e variância
$\sigma_1^2 + \sigma_2^2$.
\end{corollary}

\begin{proof}
Pela independência, o \Cref{thm:fc-soma-independentes} dá
$\varphi_{X_1+X_2} = \varphi_{X_1}\varphi_{X_2}$, e o \Cref{ex:normal}
fornece $\varphi_{X_i}(t) = \exp(-\sigma_i^2 t^2/2)$. Logo, usando
$e^{a}e^{b} = e^{a+b}$,
\[
  \varphi_{X_1+X_2}(t)
  = \exp\Bigl(-\frac{\sigma_1^2 t^2}{2}\Bigr)
    \exp\Bigl(-\frac{\sigma_2^2 t^2}{2}\Bigr)
  = \exp\Bigl(-\frac{(\sigma_1^2+\sigma_2^2)\,t^2}{2}\Bigr).
\]
Pelo \Cref{ex:normal}, o membro direito é a função característica de uma
variável $Y$ normal de média $0$ e variância $\sigma_1^2+\sigma_2^2$; como
$\varphi_{X_1+X_2}=\varphi_Y$, a unicidade do \Cref{thm:inversao} dá $X_1+X_2\overset{d}{=}Y$.
\end{proof}

\medskip

O argumento do \Cref{cor:soma-normais} é o protótipo do uso das funções
características: diretamente, seria preciso calcular a convolução
\[
  f_{X_1+X_2}(z) = \int_{-\infty}^{\infty}
  \frac{1}{\sqrt{2\pi\sigma_1^2}}
  \exp\Bigl(-\frac{(z-y)^2}{2\sigma_1^2}\Bigr)
  \frac{1}{\sqrt{2\pi\sigma_2^2}}
  \exp\Bigl(-\frac{y^2}{2\sigma_2^2}\Bigr)\,dy
\]
e completar o quadrado no expoente. O caminho pelas funções características
troca essa integral por uma multiplicação de exponenciais: a independência vira
produto e a unicidade da inversão converte identidades entre funções em
identidades entre distribuições. Esse é o mecanismo dos teoremas limites.

\subsection{O caso integrável}\label{subsec:caso-integravel}

A fórmula de inversão fica bem mais simples quando $\varphi$ é integrável:
nesse caso o limite em $T$ pode ser passado para dentro da integral, e a
medida passa a ser descrita por uma integral sem limite. Veremos que isso só
acontece quando a medida subjacente é regular, pois a integrabilidade de
$\varphi$ é uma hipótese forte e obriga $\mu$ a ter densidade.

\begin{theorem}[Densidade quando $\varphi$ é integrável]\label{thm:densidade-fc-integravel}
Seja $\mu$ uma probabilidade em $\mathbb{R}$, com função característica $\varphi$. Se
\[
\int_{-\infty}^{\infty} |\varphi(t)|\,dt<\infty,
\]
então $\mu$ tem densidade $f$ contínua e limitada, dada por
\[
f(y)=\frac{1}{2\pi}\int_{-\infty}^{\infty} e^{-ity}\varphi(t)\,dt,
\qquad y\in\mathbb{R}.
\]
Mais precisamente, para todos $x\in\mathbb{R}$ e $h>0$ vale
$\mu\big((x,x+h]\big)=\int_{x}^{x+h} f(y)\,dy$.
\end{theorem}

\begin{proof}
A ideia vale a pena ser isolada antes das contas. Na fórmula de inversão,
\Cref{thm:inversao}, aparece o núcleo $\frac{e^{-ita}-e^{-itb}}{it}$
multiplicado por $\varphi(t)$; esse núcleo é limitado por $b-a$, de modo que o
produto é dominado por uma função integrável em $t$, independente de $T$. Logo
a integral truncada converge absolutamente, e o limite em $T$ pode ser passado
para dentro da integral pelo teorema da convergência dominada. É exatamente
esse passo que torna a fórmula explícita.

Comecemos pela cota do núcleo. Escrevendo o quociente como uma integral e majorando em módulo,
\begin{equation}\label{eq:cota-nucleo-inversao}
\left|\frac{e^{-ita}-e^{-itb}}{it}\right|
=\left|\int_{a}^{b} e^{-ity}\,dy\right|
\leqslant |b-a|,
\qquad a<b .
\end{equation}
A primeira igualdade é imediata: como $\frac{d}{dy}e^{-ity}=-it\,e^{-ity}$, o
teorema fundamental do cálculo dá
$\int_{a}^{b} e^{-ity}\,dy=\frac{e^{-ita}-e^{-itb}}{it}$, e a desigualdade
segue de $|e^{-ity}|=1$ e do comprimento do intervalo, o que controla o núcleo
inclusive perto de $t=0$.

Suponha agora $\int|\varphi(t)|\,dt<\infty$ e tome $a<b$. Por
\eqref{eq:cota-nucleo-inversao}, o integrando
$\frac{e^{-ita}-e^{-itb}}{it}\varphi(t)$ é dominado, em módulo, por
$(b-a)|\varphi(t)|$, que é integrável e não depende de $T$. Em particular a
integral imprópria converge absolutamente, e o limite em $T$ que aparece em
\Cref{thm:inversao} pode ser passado para dentro da integral pelo teorema da
convergência dominada. Obtemos assim, sem nenhum limite residual,
\begin{equation}\label{eq:inversao-integravel}
\mu(a,b)+\tfrac{1}{2}\mu(\{a,b\})
=\frac{1}{2\pi}\int_{-\infty}^{\infty}
\frac{e^{-ita}-e^{-itb}}{it}\,\varphi(t)\,dt ,
\end{equation}
com a integral sobre toda a reta, e não apenas como limite de truncamentos
simétricos. A mesma cota produz ainda a estimativa
\[
\mu(a,b)+\tfrac{1}{2}\mu(\{a,b\})
\leqslant \frac{b-a}{2\pi}\int_{-\infty}^{\infty}|\varphi(t)|\,dt .
\]

Essa desigualdade já mostra que $\mu$ não tem massas pontuais. Com efeito,
fixe $a$ e tome uma sequência $b_n\downarrow a$. O lado esquerdo tende a
$\mu(\{a\})$, porque $\mu(a,b_n)\to\mu(\varnothing)=0$, por continuidade
monotônica da medida, e $\mu(\{a,b_n\})=\mu(\{a\})+\mu(\{b_n\})$, com
$\mu(\{b_n\})\to 0$: os pontos $b_n$ são dois a dois distintos, de modo que as
massas $\mu(\{b_n\})$ são os termos de uma série convergente, pois sua soma é
limitada por $1$. O lado direito tende a $0$, pois $b_n-a\to 0$ e a integral
de $|\varphi|$ é finita. Logo $\mu(\{a\})=0$; como $a$ é arbitrário, $\mu$ não
tem massas pontuais.

Podemos agora calcular a medida de um intervalo qualquer pela fórmula de
inversão. Como não há massas pontuais, o termo de correção em
\eqref{eq:inversao-integravel} se anula, e para $a=x$ e $b=x+h$ obtemos
\begin{align*}
\mu(x,x+h)
&=\frac{1}{2\pi}\int_{-\infty}^{\infty}
\frac{e^{-itx}-e^{-it(x+h)}}{it}\,\varphi(t)\,dt \\
&=\frac{1}{2\pi}\int_{-\infty}^{\infty}
\left(\int_{x}^{x+h} e^{-ity}\,dy\right)\varphi(t)\,dt
=\int_{x}^{x+h}
\left(\frac{1}{2\pi}\int_{-\infty}^{\infty} e^{-ity}\varphi(t)\,dt\right)dy ,
\end{align*}
onde a segunda linha usa outra vez a representação do núcleo como integral e a
última é o teorema de Fubini. A troca de ordem está justificada: o integrando
$\left|e^{-ity}\varphi(t)\right|=|\varphi(t)|$ é integrável em $t$ sobre toda
a reta, enquanto $y$ percorre o intervalo compacto $[x,x+h]$, de modo que a
integral dupla do módulo vale $h\int|\varphi(t)|\,dt<\infty$.

A última expressão diz que, definindo
\[
f(y)\equiv\frac{1}{2\pi}\int_{-\infty}^{\infty} e^{-ity}\varphi(t)\,dt,
\]
temos $\mu(x,x+h)=\int_{x}^{x+h} f(y)\,dy$ para todos $x$ e $h>0$. Como os
extremos têm massa nula, a escolha entre intervalo aberto e fechado é
indiferente, e $f$ é uma densidade de $\mu$. Resta ver que $f$ é contínua e
limitada. A limitação é imediata, pois
\[
|f(y)|\leqslant \frac{1}{2\pi}\int_{-\infty}^{\infty}|\varphi(t)|\,dt<\infty
\]
para todo $y$. Para a continuidade, tome $y_n\to y$. Então
$e^{-ity_n}\varphi(t)\to e^{-ity}\varphi(t)$ para cada $t$, e
$\left|e^{-ity_n}\varphi(t)\right|=|\varphi(t)|$ é um dominador integrável,
independente de $n$. O teorema da convergência dominada dá $f(y_n)\to f(y)$, e
a prova está completa.
\end{proof}

O resultado mostra que há um compromisso entre o comportamento de $\varphi$ no
infinito e a regularidade de $\mu$: quanto mais rapidamente $\varphi$ decai,
mais suave é a densidade $f$, já que as derivadas de $f$ correspondem a
integrais de $t^k\varphi(t)$, que convergem justamente quando $\varphi$ decai
rápido. A recíproca, porém, não vale: a distribuição de Cauchy, tratada em
\Cref{ex:cauchy}, tem densidade, mas sua função característica não é
integrável, e a medida só é recuperada por um limite em $T$, como discutimos
em \Cref{rem:inversao-limite}.

\subsection{Duas funções características obtidas por inversão}\label{subsec:exemplos-inversao}

Aplicando a fórmula de inversão do \Cref{thm:densidade-fc-integravel} às funções
características obtidas no \Cref{ex:triangular} e no \Cref{ex:exponencial-bilateral},
obtemos duas novas funções características. Nos dois casos o integrando da fórmula de
inversão é, a menos de uma constante, uma função característica já calculada na
\Cref{sec:exemplos}, e a identidade, lida na direção oposta, diz que a densidade
recuperada tem por função característica a densidade de partida.

\medskip

\begin{example}[Distribuição de Pólya]\label{ex:distribuicao-polya}
A função $f(x)=(1-\cos x)/(\pi x^{2})$, para $x\neq0$, com $f(0)=1/(2\pi)$, é uma
densidade de probabilidade em $\mathbb{R}$, e a função característica correspondente é
\[
\varphi(t)=\bigl(1-|t|\bigr)^{+}, \qquad t\in\mathbb{R},
\]
onde $a^{+}\equiv\max\{a,0\}$ denota a parte positiva de $a$.
\end{example}

\begin{proof}
No \Cref{ex:triangular} obtivemos
$\varphi_{\triangle}(s)=2(1-\cos s)/s^{2}$ para $s\neq0$, com
$\varphi_{\triangle}(0)=1$. Ela é integrável: perto da origem é contínua com limite
$1$, pois $1-\cos s\sim s^{2}/2$, e para $|s|\geqslant1$ vale
$\bigl|\varphi_{\triangle}(s)\bigr|\leqslant4/s^{2}$. As hipóteses do
\Cref{thm:densidade-fc-integravel} estão satisfeitas, e a densidade triangular admite,
para todo $y\in\mathbb{R}$, a representação
\[
f_{\triangle}(y)
=\frac{1}{2\pi}\int_{\mathbb{R}}\frac{2\,(1-\cos s)}{s^{2}}\,e^{-isy}\,ds
=\int_{\mathbb{R}}\frac{1-\cos s}{\pi s^{2}}\,e^{-isy}\,ds
=\bigl(1-|y|\bigr)^{+},
\]
em que a segunda igualdade absorve a constante e a última é o valor da densidade dado
no \Cref{ex:triangular}.

Fazendo $y=0$, segue que
$\int_{\mathbb{R}}(1-\cos s)/(\pi s^{2})\,ds=1$, pois o integrando é limitado perto de
$0$ e majorado por $2/(\pi s^{2})$ longe de $0$. A função $f$ do enunciado, que
coincide com esse integrando fora da origem, é portanto não negativa e tem integral
$1$; o valor $f(0)=1/(2\pi)$ apenas a torna contínua e não altera a integral.

Trocando agora os nomes das variáveis, com $s=x$ para a variável de integração e
$y=-t$ para o ponto de avaliação, e usando que $\bigl(1-|y|\bigr)^{+}$ é par,
\[
\int_{\mathbb{R}}\frac{1-\cos x}{\pi x^{2}}\,e^{ixt}\,dx
=\bigl(1-|t|\bigr)^{+}, \qquad t\in\mathbb{R}.
\]
O lado esquerdo é a função característica de $f$ avaliada em $t$, pela própria
definição, logo $\varphi(t)=\bigl(1-|t|\bigr)^{+}$ e, em particular, $\varphi(0)=1$.
Note que $\bigl(1-|t|\bigr)^{+}$ é contínua, vale $1$ em $t=0$ e se anula para
$|t|\geqslant1$: seu gráfico é um triângulo de base $[-1,1]$ e altura $1$, o mesmo
perfil da densidade do \Cref{ex:triangular} com o papel das variáveis trocado.
\end{proof}

A distribuição deste exemplo não tem nome consagrado na literatura. O nome que usamos
se deve ao papel que essa densidade desempenha em um critério de construção de funções
características devido a Pólya.

\medskip

\begin{example}[Distribuição de Cauchy]\label{ex:cauchy}
A função $f(x)=1/\bigl(\pi(1+x^{2})\bigr)$, $x\in\mathbb{R}$, é uma densidade de
probabilidade, e a função característica correspondente é
\[
\varphi(t)=e^{-|t|}, \qquad t\in\mathbb{R}.
\]
\end{example}

\begin{proof}
No \Cref{ex:exponencial-bilateral} obtivemos
$\varphi_{L}(s)=1/(1+s^{2})$, onde o índice lembra o nome de Laplace, que é
integrável: é contínua e limitada em $[-1,1]$ e, para $|s|\geqslant1$, vale
$\varphi_{L}(s)\leqslant1/s^{2}$. Pelo \Cref{thm:densidade-fc-integravel}, a densidade
da exponencial bilateral é dada pela integral de inversão, e multiplicar a identidade
por $2$ fornece
\begin{align*}
\frac{1}{2\pi}\int_{\mathbb{R}}\frac{1}{1+s^{2}}\,e^{-isy}\,ds
&=\tfrac12\,e^{-|y|},
\\[0.2cm]
\int_{\mathbb{R}}\frac{1}{\pi(1+s^{2})}\,e^{-isy}\,ds
&=e^{-|y|}, \qquad y\in\mathbb{R}.
\end{align*}
Fazendo $y=0$ na segunda identidade e lembrando que a primitiva de $1/(1+s^{2})$ é
$\arctan s$, o lado esquerdo vale
$\frac{1}{\pi}\bigl(\frac{\pi}{2}+\frac{\pi}{2}\bigr)=1$, cálculo de integração já
conhecido; logo o integrando é uma densidade de probabilidade. Trocando $s=x$ e
$y=-t$, e usando de novo que $e^{-|y|}$ é par,
\[
\int_{\mathbb{R}}\frac{1}{\pi(1+x^{2})}\,e^{ixt}\,dx=e^{-|t|},
\qquad t\in\mathbb{R},
\]
o que diz que $\varphi(t)=e^{-|t|}$ é a função característica da densidade de Cauchy;
em particular $\varphi(0)=1$.
\end{proof}

A distribuição de Cauchy é o exemplo clássico de distribuição sem média:
$\mathbb{E}\bigl[|X|\bigr]=\frac{2}{\pi}\int_{0}^{\infty}x/(1+x^{2})\,dx=\infty$,
pois a primitiva do integrando é $\tfrac12\log(1+x^{2})$, que diverge. Isso se reflete
em $t=0$: a função $\varphi(t)=e^{-|t|}$ é derivável em todo $t\neq0$, com
$\varphi'(t)=-\operatorname{sgn}(t)\,e^{-|t|}$, mas não é derivável em $t=0$, onde as
derivadas laterais valem $-1$, à direita, e $+1$, à esquerda. Como $\varphi>0$, o
módulo também não é derivável na origem: $|\varphi(t)|=e^{-|t|}$ tem o mesmo bico.

Em ambos os exemplos as densidades são simétricas e as funções características
resultantes são reais e pares, como previsto no \Cref{cor:fc-real}. Além disso, em
nenhum dos dois casos $\varphi$ é derivável em $t=0$, o que é coerente com a ausência
de primeiro momento finito nessas distribuições; voltaremos a essa relação na
\Cref{sec:momentos}.

\section{Convergência Fraca e o Teorema de Continuidade}\label{sec:continuidade}

Recorde, das notas de aula sobre modos de convergência, que uma sequência de
variáveis aleatórias $X_n$ converge em distribuição para $X$, e escrevemos
$X_n\Rightarrow X$, quando
$\mathbb{P}(X_n\leqslant x)\to\mathbb{P}(X\leqslant x)$ em todo ponto $x$ em
que a função de distribuição de $X$ é contínua. A mesma noção se formula
diretamente para medidas de probabilidade: dizemos que $\mu_n$ converge
fracamente para $\mu$, e escrevemos $\mu_n\Rightarrow\mu$, quando
$\int f\,d\mu_n\to\int f\,d\mu$ para toda função $f$ limitada e contínua. As
duas formulações são equivalentes, e é esta última que usaremos aqui. A
equivalência é o teorema de Portmanteau; o sentido que mais empregaremos, de que
a convergência das funções de distribuição acarreta a convergência das
integrais de funções limitadas e contínuas, prova-se aproximando essas funções
por combinações lineares de funções escada e usando que os pontos de
descontinuidade da função de distribuição limite formam um conjunto enumerável.

O passo seguinte em direção ao teorema central do limite é relacionar a
convergência das funções características com a convergência fraca das medidas.
O resultado que faz essa ponte, conhecido como teorema de continuidade, mostra
que a aplicação que associa a cada medida de probabilidade a sua função
característica é contínua nos dois sentidos, desde que o limite das funções
características seja contínuo na origem.

Antes de enunciá-lo, recordamos a noção de que precisaremos na prova. Dizemos
que uma sequência $(\mu_n)_{n\geqslant1}$ de medidas de probabilidade em
$\mathbb{R}$ é \emph{justa} (do inglês \emph{tight}) se, para todo
$\varepsilon>0$, existe $M>0$ tal que
\[
  \mu_n\bigl([-M,M]\bigr)\geqslant 1-\varepsilon
  \qquad\text{para todo } n\geqslant1.
\]
Em palavras, uma sequência é justa quando nenhuma de suas medidas escapa para o infinito:
toda a massa fica contida em um intervalo fixo, a menos de $\varepsilon$, uniformemente em $n$.

\begin{theorem}[Teorema de continuidade]\label{thm:continuidade}
Sejam $\mu_n$, $1\leqslant n\leqslant\infty$, medidas de probabilidade em
$\mathbb{R}$ e denote por $\varphi_n$ a função característica de $\mu_n$, isto
é,
\[
  \varphi_n(t)=\int e^{itx}\,\mu_n(dx),
  \qquad t\in\mathbb{R}.
\]

\noindent (i) Se $\mu_n\Rightarrow\mu_\infty$, então $\varphi_n(t)\to\varphi_\infty(t)$ para todo $t\in\mathbb{R}$.

\smallskip
\noindent (ii) Se $\varphi_n(t)$ converge pontualmente para um limite $\varphi(t)$ que é contínuo em $0$, então a sequência de distribuições $\mu_n$ é justa
e converge fracamente para a medida $\mu$ cuja função característica é $\varphi$.
\end{theorem}

\begin{proof}
O item (i) é o sentido fácil, e nele fica claro que a convergência fraca é uma
noção mais forte do que a convergência pontual das funções características. De
fato, fixemos $t\in\mathbb{R}$. A função $x\mapsto e^{itx}$ é limitada, pois
$|e^{itx}|=1$ para todo $x$, e é contínua, pois é a composta da aplicação
linear $x\mapsto tx$ com a exponencial complexa. Pela caracterização da
convergência fraca em termos de funções limitadas e contínuas,
$\mu_n\Rightarrow\mu_\infty$ implica
\[
  \int e^{itx}\,\mu_n(dx)\longrightarrow\int e^{itx}\,\mu_\infty(dx),
\]
ou seja, $\varphi_n(t)\to\varphi_\infty(t)$, como queríamos.

Passemos ao item (ii). Nosso primeiro objetivo é provar que a sequência é justa.
Começamos com um cálculo que pode parecer misterioso, mas que se revelará muito útil:
trata-se de integrar $1-e^{itx}$ em $t$ sobre um intervalo simétrico.

\smallskip
\noindent\textbf{Etapa 1 (o cálculo da cauda).}
Sejam $u>0$ e $x\in\mathbb{R}$. Separando a exponencial em parte real e parte imaginária,
\[
  \int_{-u}^{u}\bigl(1-e^{itx}\bigr)\,dt
  =2u-\int_{-u}^{u}\cos(tx)\,dt-i\int_{-u}^{u}\sin(tx)\,dt .
\]
A última integral se anula porque $t\mapsto\sin(tx)$ é uma função ímpar
integrada em um intervalo simétrico. Já para a primeira, como a primitiva de
$\cos(tx)$ em $t$ é $\sin(tx)/x$ quando $x\neq0$,
\[
  \int_{-u}^{u}\cos(tx)\,dt
  =\frac{\sin(tx)}{x}\Big|_{t=-u}^{t=u}
  =\frac{2\sin(ux)}{x},
\]
e a mesma expressão vale para $x=0$ por continuidade, já que $2\sin(ux)/x\to2u$ quando $x\to0$.
Obtemos assim
\[
  \int_{-u}^{u}\bigl(1-e^{itx}\bigr)\,dt
  =2u-\frac{2\sin(ux)}{x},
  \qquad x\in\mathbb{R}.
\]

\smallskip
\noindent\textbf{Etapa 2 (integrar em $\mu_n$).}
Dividindo a identidade da Etapa 1 por $u$ e integrando em $x$ com respeito a
$\mu_n(dx)$, obtemos
\[
  \frac{1}{u}\int_{-u}^{u}\bigl(1-\varphi_n(t)\bigr)\,dt
  =2\int\Bigl(1-\frac{\sin(ux)}{ux}\Bigr)\,\mu_n(dx),
\]
onde usamos o teorema de Fubini no lado esquerdo. O uso de Fubini é legítimo
porque o integrando $(t,x)\mapsto1-e^{itx}$ é limitado por $2$ e a medida
produto $dt\otimes\mu_n$ é finita em $[-u,u]\times\mathbb{R}$. Observe ainda
que o lado esquerdo é real, pois $\varphi_n(-t)=\overline{\varphi_n(t)}$, de
modo que a parte imaginária de $\varphi_n$ é ímpar e se cancela na integral
sobre o intervalo simétrico.

\smallskip
\noindent\textbf{Etapa 3 (uma cota para o integrando).}
Para limitar o lado direito, usamos a desigualdade elementar
$|\sin y|\leqslant|y|$ para todo $y\in\mathbb{R}$. De fato, se $y\geqslant0$,
\[
  |\sin y|=\Bigl|\int_{0}^{y}\cos z\,dz\Bigr|\leqslant\int_{0}^{y}|\cos z|\,dz\leqslant y,
\]
e para $y<0$ usamos a paridade de ambos os lados. Com $y=ux$ segue que
$\bigl|\sin(ux)/(ux)\bigr|\leqslant1$ e, em particular,
$1-\sin(ux)/(ux)\geqslant0$. Mais ainda, quando $|x|\geqslant 2/u$ temos
$|ux|\geqslant2$ e, usando $|\sin(ux)|\leqslant1$,
\[
  1-\frac{\sin(ux)}{ux}
  \ \geqslant\ 1-\frac{1}{|ux|}.
\]
Descartando a integral sobre o intervalo $(-2/u,2/u)$, onde o integrando é apenas não negativo,
e aplicando a cota anterior no resto,
\[
  2\int\Bigl(1-\frac{\sin(ux)}{ux}\Bigr)\mu_n(dx)
  \ \geqslant\ 2\int_{\{|x|\geqslant 2/u\}}\Bigl(1-\frac{1}{|ux|}\Bigr)\mu_n(dx)
  \ \geqslant\ \mu_n\bigl(\{x:|x|>2/u\}\bigr),
\]
sendo a última desigualdade consequência de $|ux|\geqslant2$, que dá
$2\bigl(1-1/|ux|\bigr)\geqslant1$.

\smallskip
\noindent\textbf{Etapa 4 (justeza).}
Como $\varphi$ é contínua em $0$ e $\varphi(0)=\lim_n\varphi_n(0)=1$, temos
$\varphi(t)\to1$ quando $t\to0$, e portanto
\[
  \frac{1}{u}\int_{-u}^{u}\bigl(1-\varphi(t)\bigr)\,dt\ \longrightarrow\ 0
  \qquad\text{quando } u\to0 .
\]
Fixado $\varepsilon>0$, escolhemos $u>0$ de modo que essa integral seja menor
que $\varepsilon$. Por outro lado, $\varphi_n(t)\to\varphi(t)$ para cada $t$ e
$|1-\varphi_n(t)|\leqslant2$, de modo que o teorema da convergência dominada,
aplicado na variável $t$ ao intervalo limitado $[-u,u]$, garante
\[
  \frac{1}{u}\int_{-u}^{u}\bigl(1-\varphi_n(t)\bigr)\,dt
  \ \longrightarrow\ \frac{1}{u}\int_{-u}^{u}\bigl(1-\varphi(t)\bigr)\,dt
  \ <\ \varepsilon .
\]
Logo existe $N$ tal que, para todo $n\geqslant N$,
\[
  \varepsilon\ \geqslant\ \frac{1}{u}\int_{-u}^{u}\bigl(1-\varphi_n(t)\bigr)\,dt
  \ \geqslant\ \mu_n\bigl(\{x:|x|>2/u\}\bigr),
\]
combinando a Etapa 2 com a cota da Etapa 3. Assim
$\mu_n([-2/u,2/u])\geqslant1-\varepsilon$ para todo $n\geqslant N$. Para os
finitos índices $n<N$, cada medida $\mu_n$ é de probabilidade, de modo que
existe $M\geqslant 2/u$ com $\mu_n([-M,M])\geqslant1-\varepsilon$ para todo
$n<N$ (aumentando $M$ se necessário, já que $\mu_n([-M,M])\to1$ quando
$M\to\infty$ ). Concluímos que a sequência $\mu_n$ é justa.

\smallskip
\noindent\textbf{Etapa 5 (identificação do limite).}
Falta mostrar que a sequência inteira converge, e não apenas uma subsequência.
Suponha que $\mu_{n(k)}\Rightarrow\mu$ para alguma subsequência. Pelo item
(i), já provado, a medida limite $\mu$ tem função característica $\varphi$; em
particular, o limite fraco de qualquer subsequência convergente é sempre a
mesma medida $\mu$. Por outro lado, a justeza garante que toda subsequência
possui uma subsubsequência que converge fracamente para alguma medida de
probabilidade. Esse é o teorema de Prokhorov no caso da reta: ele se obtém
aplicando o argumento diagonal de Cantor às funções de distribuição restritas
aos racionais e estendendo o limite por monotonicidade, e é a justeza que
impede a massa de escapar e garante que o limite seja ainda uma medida de
probabilidade. Seja agora $f$ limitada e
contínua. Aplicando o raciocínio anterior à sequência de números reais
$\int f\,d\mu_n$, obtemos que toda subsequência de
$\bigl(\int f\,d\mu_n\bigr)_{n\geqslant1}$ tem uma subsubsequência que
converge para $\int f\,d\mu$. Pelo critério de subsequências para sequências
de números reais, a sequência inteira converge para esse limite, isto é,
\[
  \int f\,d\mu_n\ \longrightarrow\ \int f\,d\mu .
\]
Como isso vale para toda função $f$ limitada e contínua, a caracterização da
convergência fraca em termos de funções limitadas e contínuas dá
$\mu_n\Rightarrow\mu$, o que conclui a prova.
\end{proof}

\begin{remark}[Por que a continuidade em $0$ é necessária]\label{rem:continuidade-origem}
A hipótese de continuidade do limite na origem não é um detalhe técnico.
Considere $\mu_n$ a distribuição normal de média $0$ e variância $n$.
A função característica correspondente é
\[
  \varphi_n(t)=e^{-nt^2/2},
  \qquad t\in\mathbb{R}.
\]
Para $t\neq0$ fixo, o expoente $-nt^2/2\to-\infty$, de modo que $\varphi_n(t)\to0$;
já para $t=0$ temos $\varphi_n(0)=1$ para todo $n$.
Portanto o limite pontual existe e é
\[
  \varphi(t)=
  \begin{cases}
    1, & t=0,\\
    0, & t\neq0,
  \end{cases}
\]
que não é contínuo em $0$.

Vejamos que, de fato, as medidas não convergem fracamente. Escrevendo $\mu_n$
como a distribuição de $\sqrt{n}\,Z$, com $Z$ normal padrão, para cada
$x\in\mathbb{R}$ fixo temos
\[
  \mu_n\bigl((-\infty,x]\bigr)
  =\mathbb{P}\bigl(Z\leqslant x/\sqrt{n}\bigr)
  \ \longrightarrow\ \mathbb{P}(Z\leqslant0)=\tfrac12,
\]
já que $x/\sqrt{n}\to0$ e a distribuição de $Z$ é contínua. Assim as funções
de distribuição convergem pontualmente para a função constante $\tfrac12$, que
não é função de distribuição de nenhuma medida de probabilidade. Pelo item (i)
do \Cref{thm:continuidade}, se houvesse convergência fraca, o limite teria
função característica contínua em toda a reta, e em particular em $0$; como o
limite pontual acima não é contínuo em $0$, a convergência fraca não pode
ocorrer. A continuidade na origem é, portanto, exatamente a hipótese que
impede a massa de "escapar para o infinito" sem que isso se manifeste no
limite pontual das funções características.
\end{remark}

\section{Momentos e Derivadas}\label{sec:momentos}

Nas notas de aula sobre funções geradoras de momentos, vimos que derivar $M_X$
na origem produz os momentos de $X$. Para a função característica, o mesmo
fenômeno ocorre com derivadas em relação a $t$, mas há uma diferença essencial:
como $|e^{itX}|=1$, as derivadas de $\varphi$ estão ligadas ao decaimento da
medida no infinito e não à existência de uma vizinhança onde a exponencial real
seja integrável. Esta seção torna essa relação precisa nos dois sentidos.

\subsection{Derivadas, momentos e o decaimento da medida}\label{subsec:momentos-derivadas}

Na prova do \Cref{thm:continuidade} obtivemos, para todo $u>0$, a cota
\begin{equation}\label{eq:cota-concentracao}
\mu\bigl\{x:|x|>2/u\bigr\}
\leqslant
u^{-1}\int_{-u}^{u}\bigl(1-\varphi(t)\bigr)\,dt .
\end{equation}
Ela diz que o comportamento de $\varphi$ perto da origem controla a cauda de
$\mu$: se $\varphi\equiv1$ em $[-u,u]$, o lado direito se anula, logo $\mu$ não
carrega massa em $\{x:|x|>2/u\}$; em geral, quanto menor for $1-\varphi$ em média
nesse intervalo, menor pode ser a massa de $\mu$ longe de $0$.

Há uma sutileza complexa em \eqref{eq:cota-concentracao}: o lado esquerdo é
real, enquanto $\varphi$ assume valores em $\mathbb{C}$. A parte imaginária não
atrapalha, porque não contribui para a integral. De fato, pelo item (b) de
\Cref{thm:propriedades-fc}, $\operatorname{Im}\varphi$ é ímpar, logo
$\int_{-u}^{u}\operatorname{Im}\varphi(t)\,dt=0$ por simetria do intervalo.
Portanto
$\int_{-u}^{u}\bigl(1-\varphi(t)\bigr)\,dt=\int_{-u}^{u}\bigl(1-\operatorname{Re}\varphi(t)\bigr)\,dt$,
que é real e não negativo, pois $\operatorname{Re}\varphi(t)\leqslant1$ pela
desigualdade $|\varphi(t)|\leqslant1$.

\begin{theorem}[Derivadas da função característica]
\label{thm:derivadas-momentos}
Se $\int|x|^n\,\mu(dx)<\infty$ para algum $n\geqslant1$, então $\varphi$ tem
derivada contínua de ordem $n$ em $\mathbb{R}$, dada por
\begin{equation*}
\varphi^{(n)}(t)=\int (ix)^n e^{itx}\,\mu(dx)
\qquad\text{para todo } t\in\mathbb{R}.
\end{equation*}
\end{theorem}

\begin{proof}
Provamos por indução sobre $k$ a afirmação de que, para cada
$k\in\{0,1,\ldots,n\}$, $\varphi$ tem derivada contínua de ordem $k$ dada por
$\varphi^{(k)}(t)=\int (ix)^k e^{itx}\,\mu(dx)$. O caso $k=0$ é a própria
definição de $\varphi$.

Suponha a afirmação válida para algum $k<n$ e escreva
$\varphi^{(k)}(t)=\int g(t,x)\,\mu(dx)$ com $g(t,x)\equiv(ix)^k e^{itx}$. Para
cada $t$ fixado o integrando é integrável, porque $|g(t,x)|=|x|^k$, enquanto a
derivada em $t$ do integrando é
$\tfrac{\partial g}{\partial t}(t,x)=(ix)^{k+1}e^{itx}$, de módulo $|x|^{k+1}$.

O ponto que faz a indução funcionar é a majoração $|x|^k\leqslant1+|x|^n$
sempre que $k\leqslant n$: de fato, se $|x|\leqslant1$ então
$|x|^k\leqslant1$, e se $|x|>1$ então $|x|^k\leqslant|x|^n$. Assim a única
hipótese global, $\int|x|^n\,\mu(dx)<\infty$, garante a integrabilidade de
$|x|^j$ para todos os expoentes $j\leqslant n$; em particular, $|x|^{k+1}$ é
uma majorante integrável, independente de $t$, para a derivada do integrando.
Podemos então aplicar o \Cref{thm:derivacao-integral} e concluir que
$\varphi^{(k)}$ é derivável, com
$(\varphi^{(k)})'(t)=\int (ix)^{k+1}e^{itx}\,\mu(dx)$.

Por fim, a continuidade dessa derivada segue da convergência dominada: o
integrando é contínuo em $t$ para cada $x$ e dominado por $|x|^{k+1}$, que é
integrável. Isso encerra o passo de indução e a prova.
\end{proof}

Se $\mathbb{E}|X|^n<\infty$, então $\varphi$ é $n$ vezes derivável em $0$ e
\[
\varphi^{(n)}(0)=\int (ix)^n\,\mu(dx)=\mathbb{E}\bigl[(iX)^n\bigr].
\]
Em particular, $\varphi'(0)=i\,\mathbb{E}[X]$ e
$\varphi''(0)=i^2\,\mathbb{E}[X^2]=-\,\mathbb{E}[X^2]$.
As derivadas de $\varphi$ na origem codificam, portanto, os momentos de $X$,
com um fator $i$ por ordem de derivação. Esse é o análogo exato da fórmula para
a função geradora de momentos; a diferença, já comentada na abertura desta
seção, é que $|e^{itx}|=1$, e nenhuma vizinhança de integrabilidade da
exponencial real é necessária.

Com as derivadas em mãos, obtemos a expansão de Taylor formal de $\varphi$ na
origem: se $\mathbb{E}|X|^n<\infty$, então
\begin{equation*}
\varphi(t)=\sum_{m=0}^{n}\frac{\mathbb{E}\bigl[(itX)^m\bigr]}{m!}+o(t^n),
\qquad\text{quando } t\to0 .
\end{equation*}
Aqui $o(t^n)$ designa uma função $g$ tal que $g(t)/t^n\to0$ quando $t\to0$, de
modo que o erro da aproximação polinomial é de ordem menor que $t^n$. A
expansão é apenas a fórmula de Taylor com resto de Peano aplicada a $\varphi$,
usando $\varphi^{(m)}(0)=\mathbb{E}[(iX)^m]$ para $m\leqslant n$.

O símbolo $o(t^n)$ não quantifica esse resto, e para os argumentos de
convergência que faremos adiante precisamos de uma cota uniforme em uma classe
de medidas. É o que a próxima subseção faz, e é o ingrediente que permitirá
demonstrar o teorema central do limite por meio de funções características.

\subsection{Uma cota para o resto de Taylor}\label{subsec:resto-taylor}

Em \Cref{thm:derivadas-momentos} vimos que as derivadas de $\varphi$ na
origem são os momentos de $X$. Para transformar essa informação em uma expansão
de Taylor, porém, não basta conhecer as derivadas: é preciso saber quão grande
é o erro cometido ao truncar a série. A cota a seguir resolve essa questão; ela
é uma afirmação puramente analítica, sobre a exponencial complexa, e por isso
vem antes de qualquer hipótese sobre $X$.

\begin{lemma}[Cota para o resto de Taylor]\label{lem:resto-taylor}
Para todo inteiro $n\geqslant0$ e todo $x\in\mathbb{R}$,
\[
\Bigl|e^{ix}-\sum_{m=0}^{n}\frac{(ix)^m}{m!}\Bigr|\leqslant\min\Bigl(\frac{|x|^{n+1}}{(n+1)!},\,\frac{2|x|^{n}}{n!}\Bigr).
\]
Quando $n=0$ a soma se reduz a $1$ e o segundo membro do mínimo deve ser lido
como $2$, de modo que a cota diz simplesmente $|e^{ix}-1|\leqslant\min(|x|,2)$.
\end{lemma}

Vale entender, antes da prova, por que o enunciado traz um mínimo de duas cotas.
O primeiro membro, $|x|^{n+1}/(n+1)!$, é a ordem de grandeza usual do resto de
Taylor de $e^{ix}$: ele descreve o erro como se descreveria o de uma série de
potências comum e é a cota adequada quando $|x|$ é pequeno. O segundo membro,
$2|x|^{n}/n!$, é pior por um fator da ordem de $|x|$ nessa região, mas cresce
como $|x|^{n}$ e não como $|x|^{n+1}$, sendo portanto o adequado para $|x|$
grande. É essa segunda cota, e não a primeira, que permitirá provar o teorema
central do limite sem supor o terceiro momento finito: ela exige apenas o
controle de $n$ momentos. Como as duas estimativas valem simultaneamente, vale
a menor delas.

\begin{proof}
A prova tem três etapas: obter uma identidade integral exata para o resto,
cotá-la de duas maneiras, cada uma eficiente em uma região, e tomar o mínimo.

\medskip
\textbf{Primeira etapa: a identidade do resto.}
Integrando por partes com $u=e^{is}$ e $dv=(x-s)^{n}ds$, de modo que
$du=ie^{is}ds$ e $v=-(x-s)^{n+1}/(n+1)$, obtemos
\[
\int_{0}^{x}(x-s)^{n}e^{is}ds=\Bigl[-\frac{(x-s)^{n+1}}{n+1}e^{is}\Bigr]_{s=0}^{s=x}+\frac{i}{n+1}\int_{0}^{x}(x-s)^{n+1}e^{is}ds .
\]
O termo de fronteira se anula em $s=x$, porque $(x-x)^{n+1}=0$, e vale
$-x^{n+1}/(n+1)$ em $s=0$; logo o colchete é $x^{n+1}/(n+1)$ e
\[
\int_{0}^{x}(x-s)^{n}e^{is}ds=\frac{x^{n+1}}{n+1}+\frac{i}{n+1}\int_{0}^{x}(x-s)^{n+1}e^{is}ds .
\]
A conta vale também para $x<0$: a integral é orientada e o teorema fundamental
do cálculo continua valendo nos extremos.

Para $n=0$ a identidade diz que $\int_{0}^{x}e^{is}ds=x+i\int_{0}^{x}(x-s)e^{is}ds$,
e o lado esquerdo é $(e^{ix}-1)/i$; multiplicando por $i$ e isolando $e^{ix}$,
\[
e^{ix}=1+ix+i^{2}\int_{0}^{x}(x-s)e^{is}ds ,
\]
que é o desenvolvimento de Taylor de ordem $1$ com resto integral: o
coeficiente de $x^{2}$ na série é $i^{2}/2!$, e a integração de $(x-s)$ produz
o $2$ que falta. Usando o caso $n=1$ e substituindo o valor da integral,
\[
e^{ix}=1+ix+\frac{i^{2}x^{2}}{2}+\frac{i^{3}}{2}\int_{0}^{x}(x-s)^{2}e^{is}ds .
\]
O padrão está claro. Se valer
\[
e^{ix}-\sum_{m=0}^{n}\frac{(ix)^{m}}{m!}=\frac{i^{n+1}}{n!}\int_{0}^{x}(x-s)^{n}e^{is}ds
\]
para um certo $n$, então a identidade aplicada ao lado direito fornece
\[
\frac{i^{n+1}}{n!}\int_{0}^{x}(x-s)^{n}e^{is}ds=\frac{i^{n+1}x^{n+1}}{(n+1)!}+\frac{i^{n+2}}{(n+1)!}\int_{0}^{x}(x-s)^{n+1}e^{is}ds ,
\]
já que $n!\,(n+1)=(n+1)!$. Isso é a mesma fórmula com $n$ trocado por $n+1$,
o que conclui a indução. Portanto, para todo $n\geqslant0$,
\[
e^{ix}-\sum_{m=0}^{n}\frac{(ix)^{m}}{m!}=\frac{i^{n+1}}{n!}\int_{0}^{x}(x-s)^{n}e^{is}ds .
\]
Toda a dificuldade está agora na integral do lado direito, que é o resto.

\medskip
\textbf{Segunda etapa: a cota boa para $|x|$ pequeno.} Como $|i^{n+1}|=1$ e
$|e^{is}|\leqslant1$, e mudando a variável para $u=|s|$ (o que também cobre
$x<0$ ),
\[
\Bigl|\frac{i^{n+1}}{n!}\int_{0}^{x}(x-s)^{n}e^{is}ds\Bigr|\leqslant\frac{1}{n!}\int_{0}^{|x|}(|x|-u)^{n}du=\frac{|x|^{n+1}}{(n+1)!} .
\]
O denominador vem daí: o $n!$ já estava na fórmula e a integral de $(x-s)^{n}$
em $(0,x)$ acrescenta o fator $n+1$. Esta é a cota do cálculo diferencial usual,
ótima para $|x|$ pequeno.

\medskip
\textbf{Terceira etapa: a cota boa para $|x|$ grande.}
Suponha $n\geqslant1$ e integre por partes outra vez, com $dv=(x-s)^{n-1}ds$:
\begin{align*}
\int_{0}^{x}(x-s)^{n-1}e^{is}ds
&=\frac{x^{n}}{n}+\frac{i}{n}\int_{0}^{x}(x-s)^{n}e^{is}ds ,
\\[0.3cm]
\frac{i}{n}\int_{0}^{x}(x-s)^{n}e^{is}ds
&=-\frac{x^{n}}{n}+\int_{0}^{x}(x-s)^{n-1}e^{is}ds .
\end{align*}
O ponto decisivo é que o termo de fronteira é ele próprio uma integral da
mesma família, $x^{n}/n=\int_{0}^{x}(x-s)^{n-1}ds$; substituindo, os dois
pedaços se fundem em uma única integral,
\[
\frac{i^{n+1}}{n!}\int_{0}^{x}(x-s)^{n}e^{is}ds=\frac{i^{n}}{(n-1)!}\int_{0}^{x}(x-s)^{n-1}\bigl(e^{is}-1\bigr)ds ,
\]
pois $\frac{i^{n}}{(n-1)!}\cdot\frac{i}{n}=\frac{i^{n+1}}{n!}$. É aqui que a
subtração de $1$ faz o seu trabalho: sem ela teríamos duas parcelas, uma do
termo de fronteira e outra da integral, e ao cotá-las separadamente voltaríamos
à ordem de grandeza $|x|^{n+1}$. Com o $-1$ dentro do parêntese resta uma só
integral, cujo integrando é uniformemente limitado, já que
$|e^{is}-1|\leqslant|e^{is}|+1\leqslant2$ para todo $s$. Essa limitação não
depende de $|x|$, e é isso que permite conservar o fator $|x|^{n}$:
\[
\Bigl|\frac{i^{n+1}}{n!}\int_{0}^{x}(x-s)^{n}e^{is}ds\Bigr|\leqslant\frac{2}{(n-1)!}\int_{0}^{|x|}(|x|-u)^{n-1}du=\frac{2|x|^{n}}{n!} .
\]

\medskip
\textbf{O caso $n=0$ e a conclusão.}
Para $n=0$ a segunda etapa já dá $|e^{ix}-1|\leqslant|x|$ e a terceira não se
aplica, pois divide por $n$; a segunda cota do mínimo é então $2$, que é a
estimativa trivial $|e^{ix}-1|\leqslant2$. Assim o enunciado vale para todo
$n\geqslant0$: para $n\geqslant1$ porque as duas últimas etapas fornecem ambas
as cotas, e vale a menor; para $n=0$ pela observação anterior.
\end{proof}

A transposição para funções características é imediata. Se
$\mathbb{E}|X|^{n}<\infty$, então
\[
\min\bigl(|tX|^{n+1},2|tX|^{n}\bigr)\leqslant2|t|^{n}|X|^{n},
\]
e o lado direito dessa desigualdade é integrável. Assim todos os termos abaixo
são finitos. Aplicando o lema a $x=tX$ e tomando esperanças, obtemos, para todo
$t\in\mathbb{R}$,
\begin{equation}\label{eq:cota-resto-taylor}
\begin{aligned}
\Bigl|\mathbb{E}\bigl[e^{itX}\bigr]
-\sum_{m=0}^{n}\mathbb{E}\Bigl[\frac{(itX)^{m}}{m!}\Bigr]\Bigr|
&\leqslant\mathbb{E}\Bigl[\Bigl|e^{itX}
-\sum_{m=0}^{n}\frac{(itX)^{m}}{m!}\Bigr|\Bigr]
\\[0.3cm]
&\leqslant\mathbb{E}\Bigl[\min\bigl(|tX|^{n+1},\,2|tX|^{n}\bigr)\Bigr].
\end{aligned}
\end{equation}
Primeiro o módulo entra na esperança, pela desigualdade triangular
(equivalentemente, pela desigualdade de Jensen aplicada a $|\cdot|$); depois se
usa a cota pontual do lema. Os denominadores foram descartados porque
$(n+1)!\geqslant1$ e $n!\geqslant1$, de modo que retirá-los só aumenta o lado
direito.

A cota \eqref{eq:cota-resto-taylor} vale para todo $t$ real, sem restrição de
tamanho, mas o interesse é o comportamento quando $t\to0$, e aí as duas
parcelas do mínimo competem: se $\mathbb{E}|X|^{n+1}<\infty$, a primeira dá
$O(|t|^{n+1})$; se dispomos apenas de $\mathbb{E}|X|^{n}<\infty$, a segunda dá
$O(|t|^{n})$. Em \Cref{thm:taylor-segunda-ordem} usaremos $n=2$: a primeira
alternativa exigiria o terceiro momento finito, e é a segunda que permite obter
a expansão de segunda ordem da função característica sob a hipótese de variância
finita apenas. Em resumo, o resto de Taylor de ordem $n$ é $O(|t|^{n+1})$ se
existirem $n+1$ momentos e $O(|t|^{n})$ se existirem apenas $n$.

\subsection{A expansão de segunda ordem e sua recíproca}\label{subsec:segunda-ordem}

O caso especial $n=2$ da cota \eqref{eq:cota-resto-taylor} será útil adiante, quando
tratarmos de somas de variáveis independentes. Com $n=2$, o polinômio de Taylor de
$\varphi$ na origem é
\[
1+it\,\mathbb{E}[X]+\frac{(it)^2}{2}\,\mathbb{E}[X^2]=1+it\,\mathbb{E}[X]-\frac{t^2}{2}\,\mathbb{E}[X^2],
\]
e o resto é majorado por $\mathbb{E}\bigl[\min\bigl(|tX|^3,2|tX|^2\bigr)\bigr]$. Como
$\mathbb{E}|X|^2<\infty$ implica $\mathbb{E}|X|<\infty$, os dois coeficientes são finitos.

\begin{theorem}[Expansão de segunda ordem]\label{thm:taylor-segunda-ordem}
Se $\mathbb{E}|X|^2<\infty$, então
\[
\varphi(t)=1+it\,\mathbb{E}[X]-\frac{t^2}{2}\,\mathbb{E}[X^2]+o(t^2),\qquad t\to0.
\]
\end{theorem}

\begin{proof}
Seja $R_2(t)\equiv\varphi(t)-1-it\,\mathbb{E}[X]+\tfrac{t^2}{2}\mathbb{E}[X^2]$ o resto, de
modo que, por \eqref{eq:cota-resto-taylor},
$|R_2(t)|\leqslant\mathbb{E}\bigl[\min\bigl(|tX|^3,2|tX|^2\bigr)\bigr]$. Vamos isolar o fator
$t^2$ dentro da esperança, para que a divisão por $t^2$ apareça explicitamente. Como
$|tx|^3=t^2|t||x|^3$ e $2|tx|^2=2t^2|x|^2$, dividindo ambos os lados por $t^2>0$ obtemos
\begin{equation}\label{eq:resto-dividido}
\frac{|R_2(t)|}{t^2}\leqslant\mathbb{E}\bigl[\min\bigl(|t|\,|X|^3,2|X|^2\bigr)\bigr].
\end{equation}
Basta mostrar que o lado direito tende a zero, o que é a afirmação $R_2(t)=o(t^2)$. Ponha
$Y_t\equiv\min\bigl(|t||X|^3,2|X|^2\bigr)$. Temos $0\leqslant Y_t\leqslant2|X|^2$, e
$2|X|^2$ é integrável por hipótese, logo $(Y_t)$ tem dominante integrável que não depende de
$t$. Para a convergência pontual, fixe $\omega$ com $|X(\omega)|<\infty$, o que vale quase
certamente porque $\mathbb{E}|X|^2$ é finito; quando $t\to0$, o primeiro argumento do mínimo
tende a $0$ e o segundo é a constante $2|X(\omega)|^2$, de sorte que $Y_t(\omega)\to0$. Pelo
teorema da convergência dominada, $\mathbb{E}[Y_t]\to0$; como $|R_2(t)|/t^2\geqslant0$, segue
de \eqref{eq:resto-dividido} que
\[
0\leqslant\limsup_{t\to0}\frac{|R_2(t)|}{t^2}\leqslant\lim_{t\to0}\mathbb{E}[Y_t]=0,
\]
isto é, $|R_2(t)|/t^2\to0$ e portanto $R_2(t)=o(t^2)$.
\end{proof}

\begin{remark}[Por que a cota com o mínimo]\label{rem:cota-minimo}
O ponto da cota \eqref{eq:cota-resto-taylor}, que envolve o mínimo de dois termos em vez do
primeiro termo que resultaria de uma aplicação ingênua da série de Taylor, é que assim a
conclusão de \Cref{thm:taylor-segunda-ordem} vale sob a hipótese $\mathbb{E}|X|^2<\infty$,
sem exigir $\mathbb{E}|X|^3<\infty$.

A cota ingênua para $n=2$, $\bigl|e^{ix}-1-ix+x^2/2\bigr|\leqslant|x|^3/3!$, majora o resto
por $\tfrac{1}{6}|t|^3\,\mathbb{E}|X|^3$ e, após a divisão por $t^2$, deixa
$\tfrac{1}{6}|t|\,\mathbb{E}|X|^3$, que só tende a zero se o terceiro momento absoluto for
finito.

O mínimo contorna essa exigência porque trata de modo diferente as duas regiões
determinadas por $|tX|$. Se $|tX|\leqslant2$, então $|tX|^3\leqslant2|tX|^2$ e o mínimo é o
termo cúbico, o termo fino, útil justamente onde $|tX|$ é pequeno; como ele carrega o fator
adicional $|t|$, decai quando $t\to0$, mesmo que $\mathbb{E}|X|^3$ seja infinito. Se
$|tX|>2$, o mínimo é $2|tX|^2$, que não envolve $|X|^3$: nessa região, onde valores grandes
de $|X|$ custariam integrabilidade, a cota é proporcional a $t^2|X|^2$, e o fator $t^2$ já
está fora da esperança, pronto para ser cancelado na divisão. Em suma, o fator extra $|t|$
da parte cúbica absorve o que seria uma potência maior de $|X|$, e a parte limitada fica
controlada por $|X|^2$.
\end{remark}

O próximo resultado mostra que a existência da segunda derivada de $\varphi$ na origem
implica a existência do segundo momento, isto é, ele é a recíproca de
\Cref{thm:taylor-segunda-ordem}. Como $\varphi$ assume valores complexos, a derivada deve
ser entendida no sentido da diferença segunda simétrica, expressão que aparece na hipótese.

\begin{theorem}[Segundo momento e diferença segunda]\label{thm:segundo-momento-finito}
Se
\[
\limsup_{h\downarrow0}\frac{\varphi(h)-2\varphi(0)+\varphi(-h)}{h^2}>-\infty,
\]
então $\mathbb{E}|X|^2<\infty$.
\end{theorem}

\begin{proof}
Seja $F$ a função de distribuição de $X$, de modo que $\varphi(t)=\int e^{itx}\,dF(x)$. Pela
linearidade da integral, aplicada ao integrando limitado,
\[
\frac{\varphi(h)-2\varphi(0)+\varphi(-h)}{h^2}=\int_{\mathbb{R}}\frac{e^{ihx}-2+e^{-ihx}}{h^2}\,dF(x).
\]
Como $e^{ihx}+e^{-ihx}=2\cos(hx)$, o integrando vale
\begin{equation}\label{eq:integrando-segunda-diferenca}
\frac{e^{ihx}-2+e^{-ihx}}{h^2}=-\frac{2\bigl(1-\cos(hx)\bigr)}{h^2}\leqslant0,
\end{equation}
pois $1-\cos(hx)\geqslant0$; e o limite fundamental $\lim_{u\to0}(1-\cos u)/u^2=1/2$, com
$u=hx$, mostra que $2\bigl(1-\cos(hx)\bigr)/h^2\to x^2$ quando $h\to0$, para cada $x$ fixado.

O sinal em \eqref{eq:integrando-segunda-diferenca} é o que autoriza usar o lema de Fatou na
forma que fornece cota inferior, pois Fatou se aplica a funções não negativas: se
$g_h\geqslant0$, então $\int\liminf_{h\to0}g_h\,dF\leqslant\liminf_{h\to0}\int g_h\,dF$.
Aplicado diretamente ao integrando negativo, o lema daria a desigualdade no sentido oposto,
inútil aqui. Tomando então $g_h(x)\equiv2\bigl(1-\cos(hx)\bigr)/h^2\geqslant0$, o simétrico
do integrando original, e lembrando que $2\varphi(0)=2$ e que
$\varphi(h)+\varphi(-h)=2\int\cos(hx)\,dF(x)$, obtemos
\begin{align*}
\int_{\mathbb{R}}x^2\,dF(x)
&= \int_{\mathbb{R}}\lim_{h\to0}\frac{2\bigl(1-\cos(hx)\bigr)}{h^2}\,dF(x)
\\[0.3cm]
&\leqslant \liminf_{h\to0}\int_{\mathbb{R}}\frac{2\bigl(1-\cos(hx)\bigr)}{h^2}\,dF(x)
\\[0.3cm]
&= \liminf_{h\to0}\frac{2\varphi(0)-\varphi(h)-\varphi(-h)}{h^2}
\\[0.3cm]
&= -\limsup_{h\to0}\frac{\varphi(h)-2\varphi(0)+\varphi(-h)}{h^2}<\infty,
\end{align*}
sendo a última desigualdade exatamente a hipótese. Como $x\mapsto x^2$ é não
negativa, a integral $\int x^2\,dF(x)$ é o segundo momento $\mathbb{E}[X^2]$, e
a cadeia mostra que esse número é finito, ou seja, que
$\mathbb{E}|X|^2<\infty$.
\end{proof}

Combinando \Cref{thm:taylor-segunda-ordem} e \Cref{thm:segundo-momento-finito}, a finitude
do segundo momento é equivalente à existência da segunda derivada de $\varphi$ na origem, no
sentido da diferença segunda simétrica. Se $\mathbb{E}|X|^2<\infty$, a expansão de segunda
ordem diz que $\bigl(\varphi(h)-2\varphi(0)+\varphi(-h)\bigr)/h^2$ converge para
$-\mathbb{E}[X^2]$; reciprocamente, a hipótese de \Cref{thm:segundo-momento-finito} garante
o segundo momento finito. Note que não se pede a existência do limite, apenas que o limite
superior não seja $-\infty$.

\section{Exercícios Propostos}\label{sec:exercicios}

Os exercícios a seguir acompanham as seções anteriores na mesma ordem em que
aparecem no texto.

\begin{exercise}[Parte real e módulo ao quadrado de uma função característica]%
\label{exercise:parte-real-modulo}
Seja $\varphi$ uma função característica. Mostre que $\operatorname{Re}\varphi$
e $|\varphi|^2$ também são funções características.

\emph{Sugestão:} em cada caso, exiba uma variável aleatória cuja função
característica seja a função pedida, trabalhando com cópias independentes de uma
variável aleatória $X$ de função característica $\varphi$ e, no primeiro caso,
com um sinal aleatório independente, usando a independência para fatorar as
esperanças.
\end{exercise}

\begin{exercise}[Inversão para massas pontuais]\label{exercise:massa-pontual}
Seja $X$ uma variável aleatória com função característica $\varphi$ e
distribuição $\mu$.
\begin{enumerate}
\item[(i)] Imitando a prova do \Cref{thm:inversao}, mostre que, para todo
$a\in\mathbb{R}$,
\[
\mu(\{a\}) = \lim_{T\to\infty}\frac{1}{2T}\int_{-T}^{T}e^{-ita}\,\varphi(t)\,dt .
\]

\item[(ii)] Suponha que $\mathbb{P}(X\in h\mathbb{Z})=1$ para algum $h>0$.
Mostre que $\varphi$ é periódica de período $2\pi/h$, isto é,
\[
\varphi\Bigl(t+\frac{2\pi}{h}\Bigr) = \varphi(t), \qquad t\in\mathbb{R},
\]
e conclua a fórmula de inversão discreta
\[
\mathbb{P}(X=x) = \frac{h}{2\pi}\int_{-\pi/h}^{\pi/h}e^{-itx}\,\varphi(t)\,dt,
\qquad x\in h\mathbb{Z}.
\]

\item[(iii)] Se $X=Y+b$ com $b\in\mathbb{R}$, então
$\mathbb{E}[e^{itX}] = e^{itb}\,\mathbb{E}[e^{itY}]$. Deduza daí que, se
$\mathbb{P}(X\in b+h\mathbb{Z})=1$, a fórmula de inversão de (ii) continua
válida, na mesma forma, para $x\in b+h\mathbb{Z}$.
\end{enumerate}

\emph{Sugestão}: em (i), troque esperança e integral como na prova do
\Cref{thm:inversao} e examine o limite do núcleo que aparece quando
$T\to\infty$. Em (ii), use que $e^{2\pi iX/h}=1$ quase certamente para obter a
periodicidade e, depois, escreva o limite de (i) como uma média de Cesàro de uma
função periódica de período $2\pi/h$, que converge para a média sobre um
período. Em (iii), aplique (ii) a $Y=X-b$ e reúna os fatores exponenciais.
\end{exercise}

\begin{exercise}[Média quadrática da função característica]%
\label{exercise:media-quadrado}
Sejam $X$ e $Y$ independentes, ambas com função característica $\varphi$ e
distribuição $\mu$. Aplicando o \Cref{exercise:massa-pontual} à variável
aleatória $X-Y$, conclua que
\[
\lim_{T\to\infty}\frac{1}{2T}\int_{-T}^{T}|\varphi(t)|^2\,dt
= \mathbb{P}(X-Y=0)
= \sum_{x}\mu(\{x\})^2 .
\]
No último passo, use a expressão de $\mathbb{P}(X=Y)$ em termos das massas
pontuais, vista nas notas de aula sobre variáveis aleatórias discretas, e a
independência de $X$ e $Y$.
\end{exercise}

\begin{remark}\label{rem:sem-massa-pontual}
A conclusão do \Cref{exercise:media-quadrado} implica que, se
$\varphi(t)\to 0$ quando $t\to\infty$, então $\mu$ não tem massas pontuais. De
fato, nesse caso o integrando tende a zero em cada ponto e é limitado por $1$,
de modo que a média em $[-T,T]$ também tende a zero; como
$\sum_{x}\mu(\{x\})^2=0$ e todas as parcelas são não negativas, cada massa
pontual se anula.

O \Cref{exercise:cantor} mostra que a recíproca é falsa: uma distribuição pode
não ter massas pontuais e, ainda assim, sua função característica não tender a
zero no infinito. No sentido oposto, o lema de Riemann--Lebesgue afirma que, se
$f\in L^1(\mathbb{R})$, então $\int_{\mathbb{R}}f(x)e^{itx}\,dx\to 0$ quando
$|t|\to\infty$; em particular, se $\mu$ tem densidade $f$, então
$\varphi(t)=\int_{\mathbb{R}}f(x)e^{itx}\,dx$ tende a zero quando
$|t|\to\infty$.
\end{remark}

\begin{exercise}[Densidade com função característica não integrável]%
\label{exercise:densidade-integral-divergente}
Dê um exemplo de medida de probabilidade $\mu$ que tenha densidade mas para a
qual
\[
\int_{-\infty}^{\infty}|\varphi(t)|\,dt = \infty .
\]
\emph{Sugestão}: dois dos exemplos anteriores, a saber, o do
\Cref{ex:uniforme} e o do \Cref{ex:exponencial}, têm essa propriedade. Em cada
caso, determine o comportamento de $|\varphi(t)|$ para $|t|$ grande, compare-o
com $1/|t|$ e conclua que a integral diverge.
\end{exercise}

\begin{exercise}[Soma de variáveis uniformes independentes]%
\label{exercise:soma-uniformes}
Suponha que $X_1,\ldots,X_n$ sejam independentes e uniformemente distribuídas no
intervalo $(-1,1)$. Mostre que, para $n\geqslant 2$, a soma
$X_1+\cdots+X_n$ tem densidade
\[
f(x) = \frac{1}{\pi}\int_{0}^{\infty}\Bigl(\frac{\sin t}{t}\Bigr)^{n}
\cos(tx)\,dt .
\]
Embora não seja evidente a partir da fórmula, $f$ é um polinômio em cada
intervalo $(k,k+1)$ com $k\in\mathbb{Z}$ e se anula fora de $[-n,n]$.

\emph{Sugestão}: obtenha a função característica da soma com o
\Cref{thm:fc-soma-independentes}, aplique a fórmula de inversão na forma do
\Cref{thm:densidade-fc-integravel}, explorando a paridade do integrando, e, para
a afirmação final, lembre que a densidade da soma é a convolução de $n$
densidades uniformes, logo é polinomial por partes em cada intervalo unitário.
\end{exercise}

\begin{exercise}[Médias de variáveis com distribuição de Cauchy]%
\label{exercise:cauchy-media}
Use o resultado do \Cref{ex:cauchy} para concluir que, se $X_1,X_2,\ldots$ são
independentes e têm distribuição de Cauchy, então, para todo $n\geqslant 1$, a
média $(X_1+\cdots+X_n)/n$ tem a mesma distribuição que $X_1$.

\emph{Sugestão}: calcule a função característica da média usando a
independência e compare-a com a função característica de $X_1$; duas variáveis
aleatórias com a mesma função característica têm a mesma distribuição, pela
fórmula de inversão (\Cref{thm:inversao}).
\end{exercise}

\begin{exercise}[Variâncias de normais convergentes]
\label{exercise:normal-variancia}
Suponha que $X_n\Rightarrow X$ e que, para cada $n\geqslant1$, $X_n$ tenha
distribuição normal de média $0$ e variância $\sigma_n^2$, com
$0\leqslant\sigma_n^2<\infty$. Mostre que existe $\sigma^2\in[0,\infty)$ tal que
$\sigma_n^2\to\sigma^2$. O limite pode ser $0$: nesse caso $X\equiv0$, a
distribuição degenerada em $0$, que interpretamos como a normal de variância
$0$. Nada se supõe de antemão sobre a convergência da sequência numérica
$(\sigma_n^2)$, nem sobre a natureza da distribuição limite.

\emph{Sugestão:} por \Cref{thm:continuidade},
$e^{-\sigma_n^2t^2/2}=\varphi_{X_n}(t)\to\varphi_X(t)$ para cada
$t\in\mathbb{R}$. Tome uma subsequência com
$\sigma_{n_k}^2\to\sigma^2\in[0,\infty]$, com a convenção $e^{-\infty}=0$; ao
longo dela, $\varphi_X(t)=e^{-\sigma^2t^2/2}$ para todo $t$. Se
$\sigma^2=\infty$, teríamos $\varphi_X(t)=0$ para todo $t\neq0$, contra
$\varphi_X(0)=1$ e a continuidade de $\varphi_X$ na origem. Logo
$\sigma^2<\infty$ e o limite não depende da subsequência escolhida; pelo
critério de subsequências para convergência de números reais,
$\sigma_n^2\to\sigma^2$.
\end{exercise}

\begin{exercise}[Somas de pares independentes convergentes]
\label{exercise:soma-independentes-convergencia}
Para cada $n\in\{1,2,\dots\}\cup\{\infty\}$, suponha que $X_n$ e $Y_n$ sejam
independentes, isto é, que o par $(X_n,Y_n)$ seja formado por variáveis
independentes; o mesmo vale para o par limite. Não se supõe independência entre
variáveis de índices diferentes. Suponha ainda que $X_n\Rightarrow X_\infty$ e
$Y_n\Rightarrow Y_\infty$. Mostre que
$X_n+Y_n\Rightarrow X_\infty+Y_\infty$.

\emph{Sugestão:} pela independência dentro de cada par e por
\Cref{thm:fc-soma-independentes},
\[
\varphi_{X_n+Y_n}(t)=\varphi_{X_n}(t)\,\varphi_{Y_n}(t),
\qquad t\in\mathbb{R},
\]
para todo $n$, inclusive $n=\infty$. Por \Cref{thm:continuidade} aplicado às
duas sequências, o produto converge para $\varphi_{X_\infty}(t)\varphi_{Y_\infty}(t)$,
que é a função característica de $X_\infty+Y_\infty$ e é contínua em $0$. Use
\Cref{thm:continuidade} na direção contrária.
\end{exercise}

\begin{exercise}[Função característica de um produto infinito]
\label{exercise:produto-infinito-qc}
Sejam $X_1,X_2,\dots$ variáveis aleatórias independentes, $S_n=X_1+\cdots+X_n$
e $\varphi_j$ a função característica de $X_j$. Suponha que $S_n\to S_\infty$
quase certamente. Mostre que
\[
  \varphi_{S_\infty}(t)=\prod_{j=1}^{\infty}\varphi_j(t),
  \qquad t\in\mathbb{R},
\]
onde o produto infinito é, por definição, o limite dos produtos parciais
$\prod_{j=1}^{n}\varphi_j(t)$, cuja existência é parte da conclusão. Nenhuma
hipótese de integrabilidade é feita sobre as $X_j$.

\emph{Sugestão:} a convergência quase certa implica convergência em
distribuição, como vimos nas notas de aula sobre modos de convergência; logo
$S_n\Rightarrow S_\infty$ e, por \Cref{thm:continuidade},
$\varphi_{S_n}(t)\to\varphi_{S_\infty}(t)$ para todo $t$. Pela independência,
uma indução sobre $n$ a partir de \Cref{thm:fc-soma-independentes} dá
$\varphi_{S_n}(t)=\varphi_1(t)\cdots\varphi_n(t)$, e basta identificar o limite.
\end{exercise}

\begin{exercise}[O produto de cossenos]
\label{exercise:produto-cossenos}
Iterando a identidade $\sin t=2\sin(t/2)\cos(t/2)$ obtemos
\[
  \frac{\sin t}{t}=\prod_{m=1}^{\infty}\cos\!\left(\frac{t}{2^{m}}\right),
  \qquad t\neq0.
\]
Prove essa identidade interpretando cada um dos lados como função característica
de uma mesma distribuição. Para isso: (i) mostre por indução que
\[
\sin t=2^{n}\sin\bigl(t/2^{n}\bigr)\prod_{m=1}^{n}\cos\bigl(t/2^{m}\bigr)
\]
e conclua que $(\sin t)/t$ é o limite dos produtos parciais, já que
$2^{n}\sin(t/2^{n})\to t$ quando $n\to\infty$; (ii) verifique que
$t\mapsto(\sin t)/t$, com o valor $1$ em $0$, é a função característica da
distribuição uniforme em $[-1,1]$; (iii) verifique que $t\mapsto\cos(t/2^{m})$
é a função característica de $Y_m$, que assume os valores $\pm2^{-m}$ com
probabilidade $1/2$ cada; (iv) aplique \Cref{exercise:produto-infinito-qc} às
variáveis independentes $Y_1,Y_2,\dots$ e mostre que $\sum_{m\geqslant1}Y_m$ é
uniforme em $[-1,1]$. No último passo, note que a série converge absolutamente,
pois é majorada por $\sum_{m\geqslant1}2^{-m}=1$, e que, escrevendo
$Y_m=2^{-m}(2U_m-1)$ com $U_m$ de Bernoulli simétrica, a soma coincide em
distribuição com $2U-1$, onde $U=\sum_{m\geqslant1}U_m2^{-m}$ é uniforme em
$[0,1]$, por ser a série dos dígitos da expansão binária. A unicidade da função
característica, dada por \Cref{thm:inversao}, fecha o argumento.
\end{exercise}

\begin{exercise}[A distribuição de Cantor]
\label{exercise:cantor}
Sejam $X_1,X_2,\dots$ independentes, cada uma assumindo os valores $0$ e $1$
com probabilidade $1/2$. Defina
\[
  X=2\sum_{j\geqslant1}\frac{X_j}{3^{j}}.
\]
A série converge quase certamente, e mesmo absolutamente, pois a soma dos
módulos de seus termos é limitada por $1$. Logo $X$ está bem definida, toma
valores em $[0,1]$ e tem a distribuição de Cantor, suportada no conjunto
ternário obtido removendo os terços médios. Calcule a função característica
$\varphi$ de $X$ e conclua que
\[
  \varphi(t)=e^{it/2}\prod_{j=1}^{\infty}\cos\!\left(\frac{t}{3^{j}}\right),
  \qquad t\in\mathbb{R}.
\]
Mostre em seguida, a partir da autossemelhança
$X\overset{d}{=}(2X_1+X')/3$, em que $X'$ é independente de $X_1$ e tem a mesma
distribuição de $X$, que $\varphi(3t)=e^{it}\cos t\,\varphi(t)$, e conclua que
$\varphi$ assume o mesmo valor nos pontos $t=3^{k}\pi$, isto é,
$\varphi(3^{k}\pi)=\varphi(\pi)$ para todo $k=0,1,2,\dots$. Interprete esse fato
à luz do decaimento das funções características de distribuições com densidade.

\emph{Sugestão:} cada parcela $2X_j/3^{j}$ assume $0$ e $2/3^{j}$ com
probabilidade $1/2$ e tem função característica
$\tfrac12\bigl(1+e^{2it/3^{j}}\bigr)=e^{it/3^{j}}\cos(t/3^{j})$; use
\Cref{exercise:produto-infinito-qc} para o produto infinito e
$\sum_{j\geqslant1}3^{-j}=1/2$ para recolher as exponenciais. Na recorrência,
note que $e^{i3^{k}\pi}\cos(3^{k}\pi)=1$ para todo $k\geqslant1$.
\end{exercise}

\begin{exercise}[Momentos da normal padrão]
\label{exercise:momentos-normal}
Seja $X$ uma variável aleatória com distribuição normal padrão. Use
\Cref{thm:derivadas-momentos} e o desenvolvimento em série de potências de
$e^{-t^{2}/2}$ para mostrar que, para todo $n\geqslant1$,
\[
  \mathbb{E}\bigl[X^{2n}\bigr]=\frac{(2n)!}{2^{n}\,n!}
  =(2n-1)(2n-3)\cdots3\cdot1\equiv(2n-1)!!,
\]
onde $(2n-1)!!$ é, por definição, o produto dos inteiros positivos ímpares que
não excedem $2n-1$, com a convenção $(-1)!!=1$. Em particular,
$\mathbb{E}[X^{2}]=1$ e $\mathbb{E}[X^{4}]=3$. Mostre também que os momentos de
ordem ímpar se anulam, por simetria da distribuição:
$\mathbb{E}[X^{2n-1}]=0$ para $n\geqslant1$.

\emph{Sugestão:} a função característica da normal padrão é
$\varphi(t)=e^{-t^{2}/2}=\sum_{n\geqslant0}\frac{(-1)^{n}}{2^{n}n!}t^{2n}$, e
\Cref{thm:derivadas-momentos} dá
$\varphi^{(2n)}(0)=i^{2n}\mathbb{E}[X^{2n}]=(-1)^{n}\mathbb{E}[X^{2n}]$;
compare os coeficientes das duas expressões de $\varphi$. Para os momentos
ímpares, use $X\overset{d}{=}-X$.
\end{exercise}

\begin{exercise}[Equicontinuidade das funções características]\label{exercise:equicontinuidade}
Dizemos que uma família de medidas de probabilidade $\{\mu_i, i\in I\}$ é \emph{justa} se
\[
  \sup_{i\in I}\mu_i\bigl([-M,M]^c\bigr)\longrightarrow0 \qquad\text{quando } M\to\infty .
\]
Denotando por $\varphi_i$ a função característica de $\mu_i$, dizemos que a
família $\{\varphi_i, i\in I\}$ é \emph{equicontínua} se, para todo
$\varepsilon>0$, existir $\delta>0$ tal que $|h|<\delta$ implica
$\bigl|\varphi_i(t+h)-\varphi_i(t)\bigr|<\varepsilon$ para todos $i\in I$ e
$t\in\mathbb{R}$.

\begin{enumerate}
  \item[(i)] Suponha que $\{\mu_i, i\in I\}$ seja justa e mostre que $\{\varphi_i, i\in I\}$ é equicontínua. Pelo item~(d) do \Cref{thm:propriedades-fc},
$\bigl|\varphi_i(t+h)-\varphi_i(t)\bigr|\leqslant\int|e^{ihy}-1|\,d\mu_i(y)$;
a cota do resto de Taylor \eqref{eq:cota-resto-taylor} com $n=0$, isto é,
$|e^{i\theta}-1|\leqslant|\theta|\wedge2$, aplicada a $\theta=hy$, dá então
  \[
    \bigl|\varphi_i(t+h)-\varphi_i(t)\bigr|\leqslant\int\bigl(|hy|\wedge2\bigr)\,d\mu_i(y).
  \]
Decomponha a integral em $\{|y|\leqslant M\}$, onde $|hy|\wedge2\leqslant|h|M$,
e em $\{|y|>M\}$, onde $|hy|\wedge2\leqslant2$; o segundo pedaço fica limitado
por $2\mu_i([-M,M]^c)$, uniformemente em $i$ pela justeza. Escolha $M$ e, em
seguida, $\delta$ em função de $\varepsilon$.

  \item[(ii)] Suponha que $\mu_n\Rightarrow\mu_\infty$ e denote por $\varphi_\infty$ a função característica de $\mu_\infty$. A família
$\{\mu_n: n\in\mathbb{N}\}\cup\{\mu_\infty\}$ é justa (uma sequência
fracamente convergente é justa, e o limite também), logo o item~(i) se aplica
a ela; além disso, pelo \Cref{thm:continuidade},
$\varphi_n(t)\to\varphi_\infty(t)$ para todo $t$. Conclua que
$\varphi_n\to\varphi_\infty$ uniformemente sobre os compactos: dado um
compacto, cubra-o por um número finito de intervalos de raio $\delta$
centrados em $t_1,\dots,t_k$ e, para $t$ no intervalo de centro $t_j$, escreva
  \[
    \bigl|\varphi_n(t)-\varphi_\infty(t)\bigr|\leqslant\bigl|\varphi_n(t)-\varphi_n(t_j)\bigr|+\bigl|\varphi_n(t_j)-\varphi_\infty(t_j)\bigr|+\bigl|\varphi_\infty(t_j)-\varphi_\infty(t)\bigr|,
  \]
usando a equicontinuidade na primeira e na terceira parcelas e a convergência
pontual nos pontos $t_j$ na parcela do meio. Argumente diretamente, sem
invocar o teorema de Arzelà--Ascoli.

  \item[(iii)] Dê um exemplo em que $\mu_n\Rightarrow\mu_\infty$, mas $\varphi_n$ não converge para $\varphi_\infty$ uniformemente em $\mathbb{R}$.
(Sugestão: tome para $\mu_n$ a massa unitária concentrada no ponto $1/n$ e
note que $\sup_{t\in\mathbb{R}}|e^{it/n}-1|=2$.)
\end{enumerate}
\end{exercise}

\begin{exercise}[Lei fraca e existência de $\varphi'(0)$]\label{exercise:pitman}
Sejam $X_1,X_2,\dots$ independentes e identicamente distribuídas, com função
característica $\varphi$, e $S_n=X_1+\cdots+X_n$.
\begin{enumerate}
\item[(i)] Se $\varphi'(0)=ia$, com $a\in\mathbb{R}$, mostre que
$S_n/n\to a$ em probabilidade. Como $\varphi_{S_n/n}(t)=\varphi(t/n)^n$ e, para
$t$ fixo,
\[
n\bigl(\varphi(t/n)-1\bigr)
=t\cdot\frac{\varphi(t/n)-1}{t/n}\longrightarrow iat,
\]
donde, com $\epsilon_n\to0$,
\[
\varphi(t/n)^n
=\Bigl(1+\frac{iat+\epsilon_n}{n}\Bigr)^{\!n}\longrightarrow e^{iat}.
\]
O \Cref{thm:continuidade} fornece então $S_n/n\Rightarrow a$, e a convergência
em distribuição para a constante $a$ equivale à convergência em probabilidade.

\item[(ii)] Reciprocamente, se $S_n/n\to a$ em probabilidade, mostre que
$\varphi(t/n)^n\to e^{iat}$ quando $n\to\infty$ pelos inteiros, com $t$ fixo.
Basta notar que $e^{itS_n/n}\to e^{ita}$ em probabilidade e que
$\bigl|e^{itS_n/n}\bigr|\leqslant1$, de modo que a convergência também se dá em
$L^1$.

\item[(iii)] Use o item~(ii) e a continuidade uniforme do item~(d) do
\Cref{thm:propriedades-fc} para mostrar que $(\varphi(h)-1)/h\to ia$ quando
$h\to0$ pelos reais positivos.
\emph{Sugestão:} fixe $t>0$ e, para $h>0$ pequeno, ponha
$n=\lfloor t/h\rfloor$, de modo que $nh\to t$ e $n\,|h-t/n|\to0$. Pelo item~(ii),
$\varphi(t/n)^n\to e^{iat}$; pela continuidade uniforme, $\varphi(h)^n$ tem o
mesmo limite, pois as duas potências diferem por no máximo $n$ vezes o módulo do
incremento de $\varphi$ em $|h-t/n|$. Logo $n(\varphi(h)-1)\to iat$ e, dividindo
por $nh\to t$, obtém-se $(\varphi(h)-1)/h\to ia$.
\end{enumerate}
Conclua que a lei fraca dos grandes números vale se, e somente se, $\varphi'(0)$
existe, e que nesse caso $\varphi'(0)=ia$.\footnote{%
O resultado deve-se a E.~J.~G.~Pitman (1956); a necessidade de demonstrar o
item~(iii) foi observada por John Walsh.}
\end{exercise}

\begin{remark}
Combinando o exercício anterior com o exemplo clássico de uma variável
aleatória com média infinita cuja função característica tem derivada na
origem, vê-se que a existência de $\varphi'(0)$ não garante
$\mathbb{E}|X|<\infty$. Um exemplo é uma variável aleatória simétrica com
caudas do tipo $c/|x|^2$, isto é, $\mathbb{P}(|X|>x)\sim c\,x^{-2}$ para $x$
grande: sua média diverge, pois
$\mathbb{E}|X|=\int_0^\infty\mathbb{P}(|X|>x)\,dx=\infty$, e no entanto
$\varphi(t)-1=O\bigl(t^2\log(1/|t|)\bigr)$ quando $t\to0$, de modo que
$(\varphi(t)-1)/t\to0$ e $\varphi'(0)=0$ existe.
\end{remark}

\begin{exercise}[Identidade para o primeiro momento]\label{exercise:primeiro-momento}
Mostre que
\[
  2\int_0^\infty\frac{1-\operatorname{Re}\varphi(t)}{\pi t^2}\,dt=\int|y|\,dF(y),
\]
com a igualdade entendida em $[0,+\infty]$. Sugestão: escreva
$1-\operatorname{Re}\varphi(t)=\int\bigl(1-\cos(ty)\bigr)\,dF(y)$ e recorde,
do \Cref{ex:distribuicao-polya}, que a densidade de Pólya
$\frac{1-\cos x}{\pi x^2}$ integra $1$ em $\mathbb{R}$; sendo ela par,
\[
  \int_0^\infty\frac{1-\cos x}{\pi x^2}\,dx=\frac12 .
\]
Com a mudança de variáveis $x=|y|t$ em
$\int_0^\infty\frac{1-\cos(ty)}{\pi t^2}\,dt$ obtém-se $\frac{|y|}{2}$, e o
teorema de Tonelli, já que o integrando é não negativo, conclui a troca das
integrais.
\end{exercise}

\begin{exercise}[Segunda derivada nula na origem]\label{exercise:segunda-derivada-nula}
Suponha que
\[
  \lim_{t\downarrow0}\frac{\varphi(t)-1}{t^2}=c>-\infty .
\]
Mostre que $\mathbb{E}[X]=0$ e que $\mathbb{E}|X|^2=-2c<\infty$. Em
particular, se $\varphi(t)=1+o(t^2)$ quando $t\to0$, então $\varphi(t)\equiv1$
e $X=0$ quase certamente.

\emph{Sugestão:} da hipótese segue que
$\bigl(1-\operatorname{Re}\varphi(t)\bigr)/t^2\leqslant C$ para
$t\in(0,\delta]$. Ora,
\[
\frac{1-\operatorname{Re}\varphi(t)}{t^2}
=\int_{\mathbb{R}}\frac{1-\cos(tx)}{t^2}\,dF(x),
\]
e o integrando é não negativo e converge para $x^2/2$ quando $t\downarrow0$,
para cada $x$ fixado. Aplicando o lema de Fatou ao longo de uma sequência
$t_k\downarrow0$ e depois o teorema da convergência monótona, obtemos
$\mathbb{E}[X^2]\leqslant2C<\infty$. Com o segundo momento finito, a expansão
de segunda ordem de \Cref{thm:taylor-segunda-ordem} dá
$\varphi(t)=1+it\,\mathbb{E}[X]-\frac{t^2}{2}\mathbb{E}[X^2]+o(t^2)$; divida
por $t^2$ e faça $t\downarrow0$ para concluir que $\mathbb{E}[X]=0$ e que
$c=-\frac{1}{2}\mathbb{E}[X^2]$.
\end{exercise}

\begin{exercise}[Convergência para zero]\label{exercise:convergencia-para-zero}
Sejam $Y_1,Y_2,\dots$ variáveis aleatórias com funções características
$\varphi_n$. Mostre que $Y_n\Rightarrow0$ se, e somente se, existe $\delta>0$
tal que $\varphi_n(t)\to1$ para todo $t$ com $|t|\leqslant\delta$. Sugestão:
numa das direções use o \Cref{thm:continuidade}; na outra, suponha a
convergência apenas em $[-\delta,\delta]$ e aplique a cota de concentração
\eqref{eq:cota-concentracao}, que dá, com $\eta=2/\delta$,
\[
  \mathbb{P}(|Y_n|>\eta)\leqslant\frac{1}{\delta}\int_{-\delta}^{\delta}\bigl(1-\operatorname{Re}\varphi_n(t)\bigr)\,dt\longrightarrow0
  \qquad\text{quando }n\to\infty .
\]
Logo $Y_n\to0$ em probabilidade, e portanto em distribuição.
\end{exercise}

\begin{exercise}[Séries que convergem em distribuição]\label{exercise:serie-converge-dist}
Sejam $X_1,X_2,\dots$ independentes e $S_n=\sum_{m\leqslant n}X_m$. Mostre
que, se $(S_n)$ converge em distribuição, então $(S_n)$ converge em
probabilidade. Conclua, a partir do teorema de equivalência de P.~Lévy para
somas de variáveis independentes, que $(S_n)$ converge quase certamente.

\emph{Sugestão:} se $\varphi_k$ é a função característica de $X_k$, o incremento
$S_m-S_n=\sum_{n<k\leqslant m}X_k$, para $m>n$, é soma de variáveis
independentes, logo
\[
  \varphi_{S_m-S_n}(t)=\prod_{k=n+1}^{m}\varphi_k(t)=\frac{P_m(t)}{P_n(t)},
  \qquad P_j(t)=\prod_{k=1}^{j}\varphi_k(t).
\]
Se $S_n\Rightarrow S$, então $P_j(t)\to\varphi_S(t)$ para todo $t$; como
$\varphi_S$ é contínua e $\varphi_S(0)=1$, escolha $\delta>0$ com
$|\varphi_S(t)|>\frac12$ para $|t|\leqslant\delta$. Nesse intervalo os
denominadores não se anulam para $n$ grande, e o quociente $P_m(t)/P_n(t)$
tende a $1$ quando $m,n\to\infty$, para cada $t$ fixado. Pelo
\Cref{exercise:convergencia-para-zero}, segue que $S_m-S_n\to0$ em
probabilidade quando $m,n\to\infty$, isto é, $(S_n)$ é de Cauchy em
probabilidade. Escolha então $n_1<n_2<\cdots$ com
$\mathbb{P}\bigl(|S_{n_{k+1}}-S_{n_k}|>2^{-k}\bigr)\leqslant2^{-k}$; o lema de
Borel--Cantelli garante que, quase certamente, $|S_{n_{k+1}}-S_{n_k}|\leqslant
2^{-k}$ a partir de um certo índice, de modo que $(S_{n_k})$ converge quase
certamente. Como essa sequência também converge em probabilidade e a sequência
inteira é de Cauchy em probabilidade, ela converge em probabilidade para o mesmo
limite. A última conclusão, de que a convergência é quase certa e não apenas em
probabilidade, é o teorema de equivalência de Lévy: para somas de variáveis
independentes, a convergência em probabilidade acarreta a convergência quase
certa. Sua prova se apoia na desigualdade maximal de Lévy e não será
desenvolvida nestas notas.
\end{exercise}

%----------------------------------------------------------------------
% BIBLIOGRAFIA
%----------------------------------------------------------------------
\bibliographystyle{plain}
\bibliography{../arquivos-aux/referencias}

\end{document}
